% ================================================================
%  18.5.5 Le seuil + 18.5.6 L'album de famille (原书 p.836–864)
%  源文: _work/chunks/chunk08.md, chunk091.md, chunk092.md,
%        chunk093a.md, chunk093b.md, chunk094a1.md, chunk094a2.md, chunk094b.md
% ================================================================
\subsection{门槛}   % 18.5.5. Le seuil

\notehead{172}{}

\pmark\origpage{1102}(3月22日)我原以为顶多一两天、撑死十来页,就能把这些自从10月起我就打算重新巡视的
出名的“四种运算”走完一遍。可如今一鼓作气干了三个多星期,页数排到了一百上
下——居然还没有全然写完!初稿从2月26日写到3月1日,就花了四天。它给我的,不过是一块底布,
好(终究)在上头绣出一桩“事情”来,而不只是一纸侦查结论。3月1日的第二天重读初稿,那印象可
不妙:一张没完没了的“诉苦状”,而且,除了三四位真正在行的专家(还得假定他们有耐心读完
……),只怕谁也读不懂。我明白:至少得约略交代一下所为何事,也就是说,至少得铺一个语境——
否则就白搭\fnote{835}{(*)}{整部《收获与播种》的思考里,除这里之外,我曾对“即兴”写作方式作过这种(规模确实小些)让步的,只有“那张条子——或新伦理”一节(n\textsuperscript{o}~33)与“不义——或归来的意义”那条注(n\textsuperscript{o}~75)。}。

这就不可避免地带来一些重复,同\emph{葬礼}第一部分相比——可有些情形,重复不但有益,甚至
必不可少(在数学里,一如在别处)。而且在这些情形里,人们很快发现:那些所谓的“重复”其实并
非重复,因为被“重说”的东西,实际上是\textbf{重看了一遍、以新的眼光再看},并且照明已经不同。在
为“四种运算”铺设“语境”、交代我著作的某些侧面时,我觉得自己对这份著作学到了些东西,把它
摆得更正了。这么做,我也许并没有就我自己或他人学到什么真正新的东西,可我不后悔费了好几
天的功夫,把那份“诉苦状”式的初稿这样重写。这份著作,我放入过我所能给出的最好的东西;它
值得我怀着成熟给我的那份回顾,以不同的光线重新去认识它。正当我准备原原本本记录这份著作
自我离去后(交在了可靠的人手里,对此我从不怀疑……)遭受的一切之时,好事在于:我在它身上
\textbf{停一停},在它的位置上、在构成它之美的那个统一体上,哪怕只停几页——以此再一次,向那
些我亲眼看见被践踏的东西,表达我的敬意。

\pmark\origpage{1103}可这还不算完,差得远呢!放弃了“诉苦状”的写法——那种夹着编号、指向\emph{葬礼}第一部分
那些血肉更丰满的注的写法——我明白了:这些我重新拿起的笔记,同《收获与播种》里所有其他
章节与笔记一样,必须自身可读,把它们要说的实质,完整地传达出来,独立于那些指向思考\textbf{另一
时刻}之注的引用。这又带来许多“并非重复的重复”——也就是,以新的照明,重看将近一年前、在发
现的新鲜震动里逐日记下的东西。何况当时,那么多意外的、有时难以置信的事实劈头盖脸,压根谈
不上一场稍具方法的“侦查”。那时,我不过竭力去接住劈头砸下来的东西,勉力把它们“安放”进去,
细节顾不上。我的精力,大部分被\textbf{招架}这些我发现的原形毕露之物的那种\textbf{癫狂}、那种难以置
信(就像那篇皇帝新衣的故事里一样……\fnote{836}{(*)}{见同名之注,n\textsuperscript{o}~77$'$。})所吞掉;尤其是,要承受那股突然扑面而来的暴戾、犬儒与轻蔑之气——“顶着那副光鲜的面孔……”我太熟悉的面孔;那是另一个时代的气息,我经历过它,我没有忘……

而这最近三个星期,反倒成了补完去年那场风暴式侦查的机会:把某些文本再往细里翻一翻(SGA~5,
尤其是那卷所谓的“SGA~4½”)。于是生出一串(有一阵眼看不象话地收不了尾的!)多少原原本本
的脚注;其中一些成了子注,而其中一个(预定的名字叫“那个公式”)占用了我整整四天,又裂成
另外四个\fnote{837}{(*)}{(6月1日)如今已经变成六个了……}……有几次,我觉得永远写不完了——可到头来,它还是收敛了
\fnote{838}{(*)}{(6月1日)当然,只是一场临时的“收敛”,因为“封神”那条注,最后炸成了三十来条各自独立的注、子注等等,单是它们就有足足一百五十页!}。眼下撂下的,还有十来页涂改得太厉害、须得重写的页,以及“分赃”与“封
神”两条注的脚注,留待日后补上。眼下这样就可以了!哪怕日后还要回头管\dmark{}\origpage{1104}“后勤”,我急于收场,
急于不再耽搁,把“四种运算”这一章里我看到的还有实质可说的东西说出来。

我在\emph{葬礼}里区分出两个“面向”或“层次”:它们紧紧相系,却毕竟不同。它们之间(至少在
我眼中)相当分明地隔着一道\textbf{门槛}。

一方面,是“时尚之风”的面向(有时直到那股“嘲弄之气”——我在《收获与播种》里已不止一次
有机会说到它)。它主要表现在我在别处\fnote{839}{(*)}{见注“掘墓人——或教团全体”,n\textsuperscript{o}~97。}称作“自动排斥的态度”的东西上:这种态度常常拦腰截断数
学常识的简单反射,冲着某些人及其数学贡献而去——这里说的是我,以及另外几个人;他们被归
类(尽管当事人有时竭力想同我划清界限)为同我“串通一气”。在我的情形,要“排斥”(或“埋葬”)
我带来的\textbf{一切},当然不可能:好大一部分,早在我1970年告别数学舞台之前,就已进入日常通用
的公共领域\fnote{840}{(*)}{不过,确实,连某些早已进入“日常”通用(至少在我学生与近旁合作者的小圈子里)的观念与技术,自我离去后也被埋葬了。不妨说,尤其如此的是$\ell$-进上同调那件工具:我在 SGA~5 里(以 SGA~4 的关键结果为起点)极详尽地发展过它。它被我那班上同调学生压在斗底下,Deligne首当其冲,直到1977年,才以我所知的那个形状与精神被掘出土来。}。不过,
确实(而且我在一年前的注“我的孤儿们”(注 n\textsuperscript{o}~46)里第一次确认了这一点):我的著作——
写下的与未写下的——在\textbf{上同调主题}上,绝大部分自我离去的次日起就被埋葬了,头一份经手
的正是我的学生。(我引入的某些主题,四年后、七年后或十二年后被掘出土来,不提我这个人——
但那我们已经触到“第二层次”了……)

这样的自动排斥,人们尽可以惋惜:它有时违背对他人的最起码的体贴与敬意,而且在任何情形下
都与数学的常识与鉴别力无关。而当它打击的是那些手段有时十分出色的年轻数学家时,就更值得
惋惜:“轻蔑之齿”咬灭了一种欢乐,糟蹋了一份原本美好的激情,让它化为苦涩——投入眼看着被
糟蹋了(按那些说一不二的共识……)。还值得惋惜的是,\dmark{}\origpage{1105}当这种排斥打击的是那些简单而丰
饶的观念——它们早已充分证明自己,从无中生有地催生出如今“人人”看也不看就使用的强大工
具。在前一种情形(一个被毁掉的志向),伤害多半不可逆转;在后一种情形则不然——
因为迟早,那些简单而根本的观念,那些“在路上”的观念,终会显形或重新显形,终会进入公共遗
产。不管怎样,谁也不能合情合理地强求任何人,对一个人、一部著作、一个观念\textbf{往好里}想——
只要他(出于只有他自己知道的理由)一心想\textbf{往坏里}想,或干脆忘了它。这类问题,当然属于“个
人伦理”,而且属于其中精细而要紧的部分;但依我看,把它弄成一桩集体的“科学伦理”问题,恐怕
不妥;倘若硬要试,只怕药比病还凶……

我方才影射的第二个“面向”或“层次”,恰恰是这种集体伦理被破坏的层次。我说的那道\textbf{门槛},
是一条\textbf{共识}——据我所知,自有科学见诸文字记载以来,它在一切科学里都得到普遍接受。这条
共识规定:谁也不得把他从\dmark{}\origpage{1106}别人那里取来的观念\fnote{841}{(*)}{我在这里说“观念”,当然,在数学里绝不单指“结果”。常常,一个提得好的\textbf{问题},它触及一个此前无人看出的关键点,就比一个“结果”更重要,哪怕更难。即便这个问题还没有凝结成一个精确的\textbf{命题}——它本可以构成一个假设回答的胚胎,或者一个(尚属猜测的)多少完整的回答——情形也是如此。当然,从一个起初含混的问题里提炼出这样一个命题,是数学工作中本质而具创造性的一部分。把一个(也许深刻的)问题加工后的版本端出来,却隐去它的出处(即便加工是这位魔术师自己做的),同隐去一个形状完整的深刻命题的出处——借口是自己给出了证明——一样,都是剽窃,正如把别人给的证明当成自己的拿出来一样。
对丰饶\textbf{观念}的引入,道理相同——它们往往比好的命题更关键:因为“好命题”的问题,只有在已经提炼出好观念之后才谈得上。这里的道理相同:借口修改过、甚至改进过取自他人的观念,就隐去其出处,其不老实,同原封不动地“借用”这观念,毫无二致。最常见的情形是:提出一个问题(哪怕还含混),拿出一个命题或一个观念(哪怕不完善、临时),才是关键的一步,而不是人们后来加上的那些改进(更精确、更广、更深)。可即便情形不是这样,对那个在改进所受之物上做出原创工作的人,这仍旧不能当作隐匿其所受之物的“理由”(或者,同一回事:把它“拆掉”……)。
如我在别处已强调过的(在注“肌肉与肚肠(阳葬阴(1))”(n\textsuperscript{o}~106)的子注 n\textsuperscript{o}~106$^{1}$里),一条猜测性命题的“价值”,既不取决于其设想的难度,也不取决于它多“可信”,也不取决于它最终被证为真或为假。无论如何,人们愿意给一个数学观念(它表达于一个问题、一个命题、一个观念或一个证明之中)或一组观念的“价值”,很大程度上是主观的,难以成为科学伦理之共识的对象。所以,一个诚实的科学家,会标明他所使用的(显式地或默会地)、不属于“众所周知”的一切观念的出处,而不肯顺从那条斜坡:对他已在内心决断(也许是为了某桩可疑的案子……)反正“显而易见”、“不足道”、“不要紧”(或同类评语)的某个观念,隐去其出处。}

\pmark 这条共识责令我们:凡我们所陈述、使用或发展的观念,每次——至少当这些观念并非出自
我们自己的手笔、也还不是公共遗产、尚未被(不是三四个入室弟子,而是)“\textbf{所有人}”知晓的
时候——都要标明来历。

我不记得曾听人质疑过这条共识。在我身属数学圈子的岁月里——从1948年(二十岁的年轻新手,
去高等师范学院听Cartan的课)到1970年(我告别数学舞台)——我极少有机会目睹、甚至极少听说
有人公然违背这条共识或原则;只在一位同事兼朋友那里,他在这一章上稍欠检点\fnote{842}{(*)}{这位同事的情形,顺带谈过,在《收获与播种》第一部分,就在紧接着下一句所引的那一节里。经过一年多的回顾,这个“案子”如今所显出的分量,是我从前没有给它的。}。正如我在《收获与播种》第一部分里强调过的
(在“一个严守的波利钦内尔秘密”一节,n\textsuperscript{o}~21),遵守这条原则,对一个具备起码的诚实与自尊
的人,绝非自然而然之事。相反,它要求极大的警觉:童年便根深蒂固的反射,自然而然地推着我们
高估自己的功劳,把我们在自己之内消化他人观念的工作,同这些观念本身的构思混为一谈——两
者实属完全不同的层次。写上面所引那一节——一年多了——的时候,我显然对自己也没弄清:该给
这条共识多大的分量。那时我心里有一片含混(在思考的这个阶段,我自己也没有清楚意识到\dmark{}\origpage{1107}),连着一种弥漫的感受:对\textbf{他人}(比方说对我自己的学生)提出严格要求——在他们
同我的关系上遵守这条原则——是小气,是一种配不上我的卑微。所以那时,我之内有一道\textbf{暧昧};
我只是在6月1日那条同名之注(n\textsuperscript{o}~63$''$)的思考里,才把它看清。这场思考把这道暧昧驱散净尽;
而我当时意识到,这道暧昧,从最初(六十年代初)直到去年,一直沉沉地压着我同学生的关系。我
懂得了:数学家这一行(或更一般地,科学家这一行)的严谨,首先意味着对自己怀有极大的警觉,
遵守这条众共识中之首要者;同时也意味着对他人同样严格的要求,尤其对那些我们负有领入门墙
之责的人。

年复一年,我越来越懂得:这一行\textbf{绝不只是}一门技术上的窍门,甚至不只是发挥想象、解
决出了名的难题的能力。从某种意义上,这一点我从来就知道——可我低估了“伦理”的一面,或者
说\textbf{集体}的一面\fnote{843}{(*)}{我在这里无意说:一个处境的“伦理”面向,总是同时就是“集体”的面向,牵涉一个人同群体(这里即“同事”或“同侪”之群)的关系。不过,在我正在检视的这条“共识”上,情形确实如此。按照自童年起塑造了我的世界观的那些特殊条件作用,直到去年,我总倾向于低估(乃至无视)集体之物,而偏重个人之物。我的“数学历险”中“集体历险”的那一面,去年才清楚地向我显现,先是在“Galois的遗产”一节(n\textsuperscript{o}~7),但尤其《收获与播种》第一部分末尾的几节:“孤独的历险”与“一段过去的重量”(n\textsuperscript{o}s~47、50)。},以为它在善意的人们、体面的人们之间,“理所当然”。就这样,我为我刚
才说的那道“暧昧”备好了土壤;那道暧昧,也是(顶着一副虚假的“慷慨”的门面)对我的学生与
类似者的\textbf{纵容},而更隐蔽的,是对我自己的一种纵容。

我离开了这个“善意的人们、体面的人们”的圈子——它也曾是\textbf{我的}世界,我曾乐于自认同属
的世界。在去年4月19日之后的几个星期里,我稍微仔细地看了它一眼:离开不到十五年,我
在里面看到一种\textbf{腐败},做梦我也想象不出它的样子。

\pmark\origpage{1108}对我而言,这是一个谜:身为这个世界的一员,“搞数学”还能有什么\textbf{意义}——如果只
是把它当作一种\textbf{权力}的手段,或者(对身份低微的)当作一种谋一碗饭的手段,物质条件嘛,
说实在的还挺舒服(只要运气好,早早就“安顿”下来了……)。

\subsection{(5) 家庭相册}   % 18.5.6. (5) L'album de famille
\notehead{173}{}

\fnote{844}{(*)}{眼下这条“家庭相册”之注,最初本是前一道注“门槛”(同日,3月22日所写)的直接续篇。这一部分,如今构成其 a.部分(“一位备受簇拥的死者”);6月10日、11日又添了 b.(“新面孔——或吊嗓”)与 c.(“众人之上的那一位——或首肯”)两部分。下一道注“升级(2)”(n\textsuperscript{o}~174)——也是3月22日——紧接本注 a.部分(同日所写)而来。至于 b.、c. 两部分的脚注,日期是6月13日与14日。最后,6月18日又添了 d.(“最后一刻——或一个禁忌的终结”)这一部分。}

\subsubsection*{a.\ 一位备受簇拥的死者}
\tocsubsub{a.\ 一位备受簇拥的死者}

\pmark (3月22日)说得粗鲁些:\emph{葬礼}里有“时尚”的层次,也有“诈骗”的层次。也许不过是
我落伍了;我那个年代被看作诈骗的行径,如今只要干它的人属于上流社会,就早已成了完全正当、
体面的行当。也许那道“门槛”,早就荡然无存了呢?“第二层次”是\textbf{一场唯一的、规模浩大的诈骗
行动}:目标是我在上同调主题上的全部著作;继它之后,又是Zoghman~Mebkhout——那位鲁莽的继承
人、被埋葬的宗师的身后学生、固执而无名之辈——的著作。这场行动的总指挥,是另一位学
生:绝非身后的,倒是隐秘的,没错;他玩着一个心照不宣的角色——我著作的“继承人”——同时对
著作、对工匠,一并拒认、一并拆毁。他是我的朋友\textbf{Pierre~Deligne}。他的干练副官,不是别
人,正是随他一道选了“上同调”这条线的四位学生:\textbf{J.L.~Verdier、L.~Illusie、P.~Berthelot、
J.P.~Jouanolou}。死者想必被簇拥得妥妥帖帖:既有同享\emph{葬礼}哀荣的\textbf{共殓诸友}
\fnote{845}{(*)}{说实话,据我所知,“共殓者”不是一位,而是\textbf{四位}——即四道棺材注(棺材一至四,n\textsuperscript{o}s~93--96)所写的那四位。},又有生前曾是他的亲近者的那些人。至于临时凑数的殡仪工——由总指挥导演
的那场双重埋葬里,前来帮衬的——我看见另外七位“世界闻名”的数学家(借某张广告招贴的\dmark{}\origpage{1109}措辞
\fnote{846}{(*)}{指 IHÉS 1983年为建所二十五周年出的纪念小册子。关于它,见“悼词(一)(二)”两条注(n\textsuperscript{o}s~104、105),尤其第454页。}),在葬礼仪式进行之中不时露面;仪式已在家庭相册(又名“四种运算”)里
巡视过一遍。他们是(按其在仪式中的分量)\textbf{Teissier、Beilinson、Bernstein、Milne、Ogus、K.Y.~Shih、Katz}
\textbf{(N.~Katz)}。

\pmark 到此,我巡视了据我所知曾以这样那样名分\textbf{积极参与}“葬礼行动”的数学家。共十二位
\fnote{847}{(*)}{与(《收获与播种》第一部分)“耶稣与十二使徒”一节里相同的“十二”——那一节巡视了所有跟我做到国家博士论文层级的学生。确实,在我葬礼的积极参与者之中——这一回,是Springer殡仪股份有限公司那一边,而非信众教团那一边——还有 K.F.~Springer博士(那家可敬公司的共同老板)与Peters、Byrne两位博士;后文(注 n\textsuperscript{o}~175)还会谈到他们。于是是十五位!}。至于后四位,就我所知的事实,我无从断定他们存心不良。但我认为,他们的责任并不因此就轻了分毫。因为,他们乐于对他们所做的事视而不见,这本身是一种选择;选择,豁免不了他们对自己行为的责任。

至于那些直接串通的参与者,要我开一张哪怕不完全的名单,或估一估他们的人数——那人数肯定
是另一个量级——我实在无能为力。只需提醒一句:这里面有1981年6月吕米尼那场“难忘讨论会”
(人称反常层会议)的全体参加者;还有那卷顶着“SGA~4½”之名的文本的读者们之中,多少知道
SGA 这个符号意味着什么、却“听之任之”的所有人。

\pmark 我看见两份文字,它们见证着七十年代、八十年代数学里的一场\textbf{凋败},在我们这门科学
的历史上,恐怕还从来没有堪与比拟的先例。其中一份,凋败在它给自己起的名字里就绽开了——那
名字本身就是一次(天才的……)冒名:那份数字标着“SGA~4½”(作为常用的参考符号)、又称
“平展上同调”的文本——P.~Deligne署名,“协作”(除 L.~Illusie与 J.L.~Verdier等人之外,还)有
A.\textbf{Grothendieck}\fnote{848}{(*)}{关于这份“协作”的含义——它是Deligne所设骗局的一部分——见注“对调”(n\textsuperscript{o}~68$'$)。}。第二份文字,是1981年6月吕米尼\dmark{}\origpage{1110}讨论会的文集;尤其并且首要的,是其第一卷:由讨
论会引论(Teissier与Verdier署名)与讨论会的主报告(Beilinson、Bernstein、Deligne署名)构成。

倘若直接参与过这场凋败的、以这样那样名分参与的人之中,有一位能在自己身上找到那份朴素与
勇气,公开道个歉——或者更好,公开就\textbf{他自己}经历的那部分原委说个明白——这对所有
人都是一桩善事,也是那一代容忍了这等凋败的数学家的光荣。可这多半是奢望。

同样是奢望的,想必还有:让 J.L.~Verdier不再在高等师范学院占着Henri~Cartan曾占的那个位置。
那无疑是法国培养数学“接班人”的关键岗位。很久以前,当我得知Verdier——他曾是我的学生,
我曾钟爱于他——升到了这个职位,我自己觉得脸上有光(同时暗地里也有几分得意)。那时,我丝毫
没有怀疑过:Verdier会完美地承担曾属于Cartan的角色,让那些最有志于数学的年轻人在与他的接触中学好这门手艺。可今天我看见(多年前已在看,却从来没有以今天这般粗暴的清晰看见)我
错了;我在这里明白地说出来,不是要给他或任何人抹黑。我认为,他已丧失了指导研究的资格。这样
说,我并不推卸我自己的一份责任:是我把他教坏了(他同我所有其他学生一样)——那门我热爱并
且依然热爱的手艺。

\subsubsection*{b.\ 新面孔——或吊嗓}
\tocsubsub{b.\ 新面孔——或吊嗓}

(6月10日)“家庭相册”这条注动笔以来,两个半月过去了。我当然没料到,还须回到它上面来——
因为对\emph{葬礼}的调查又有了新的反复。主要是那场微薄的“封神”——我当时刚写了五六页——
炸开成一场大写的、足足一百五十页密密实实的\textbf{封神},让我随之发现了一些“新面孔”:它们理应
在家庭相册里占一席之地。也有一些早已熟识的面孔,如今显出他们同样属于那支大军——那些
在“诈骗”的层次上积极参与“葬礼运算”的人。我在这里“备而勿忘”地把他们再巡视一遍\dmark{}\origpage{1111},也好让
每一位当事人确知自己与名流为伴(不过,这事想必早就办妥了……)。新到的照片,按它们
引起我注意的先后,插入相册。

先说Springer出版公司那一边:\textbf{K.F.~Springer}(公司共同老板之一)、\textbf{Peters}、与
\textbf{Byrne}夫人(C.M.~Byrne)。细节见下文“殡仪馆——‘为科学服务’”那条注(n\textsuperscript{o}~175)。写本注开头、
3月22日的那几天,我刚收到 K.F.~Springer的信(日期3月15日)没几天;那封信,把我对那家可敬
的殡仪公司——忠于“为科学服务”的箴言——的精神的最后一点疑虑,一扫而空。

封神那一边(经由对当值无名之辈的埋葬),我知晓了\textbf{Kashiwara、Hotta、Brylinski、Malgrange(B.~Malgrange)、Laumon、Remmert}诸位的贡献;还不算一位\textbf{匿名审稿人}——其存心不良,
毋庸置疑;不过说真的,倘若要把与\emph{葬礼}沾亲带故的那些可疑文章与书籍的审稿人(reviewer) 也一一
点验,那准得另备一册相册。此外,我的老朋友Katz重新露面了,这一回的语境,让我对他本来(至少)
存着的那份善意推定,化为乌有。这样一来,据我所知曾以这样那样名分积极参与那场所谓“反常讨
论会”的骗局行动的、国际知名数学家,人数到了十四位(算上那位出名的匿名审稿人,十五位)。详
情备有文件,见“封神”,尤其“……以及横财”、“光荣之日”、“黑手党”、“掠夺的空白授权——或‘高
尚行当’”几条注(n\textsuperscript{o}s~171(iii)(iv)、171$_{2}$、171$_{4}$)。

最后,在“Motif 运算”那一边,出现了(迟做总比不做好)我的又一位学生,稍稍离大队人马有点距
离。事后,我几乎是被逼着把他算作我的“上同调”学生的第六位,尽管“我那个年代”,他对上同调是
什么,连半点概念都没有。他是Neantro~Saavedra~Rivano;显然,他被人当了“棋子”(当然,出
于他自己的自愿),而非为自己行动。他与Monsieur~Verdoux(化装成“献殷勤的骑士”)\dmark{}\origpage{1112}周旋的历险,已
在“第六根钉子(钉进棺材)”那组注(n\textsuperscript{o}s~176$_{1}$至176$_{7}$)里逐页复原——那是4月19、20日写的
(最后一道除外,尚待写出)。于是,在我“从前的”学生之中,积极参与宗师\emph{葬礼}
的,人数也到了六(十二之中)。Saavedra在\emph{葬礼}中的份额,其特别之处在于:他参与的那场
“塔纳卡范畴”(原文如此)运算,是第一场大规模行动,目标是变走我著作的一个重要侧面、以及我
发展的那门哲学(就其情形而言,是在 Motif 理论的带动之下、乘其便而发展的)的父亲身份。

把新到的面孔算进相册,并撇开Springer-殡仪馆一方的贡献、只算信众教团一方的,那么据我所
知曾积极参与\emph{葬礼}——按我那个年代人们管它叫“诈骗行动”的那种水平——的知名数学家,
人数到了十九位\fnote{849}{(*)}{算上那位出名的匿名审稿人,二十位。}。这些参与者中,只有三位——即与 P.~Deligne合署“难忘文集”《数学讲义》
第900种的那三位——在我看来,存心不良尚未坐实。

这份名单,当然远没有穷尽我的同事、旧学生或朋友——他们以这样那样名分、或主动或被动地参
与了我的葬仪,却没有走到参与一场铁板钉钉的诈骗的地步。我点出过三十来位;其中大多数,已在
对\emph{葬礼}的思考中提到;算上前面那些,凑足了结结实实的五十来位——而这还只是:最近八九
年,我在遥远的隐居处,身不由己地得知的那些;或是一场刻意保持得极有限的调查之中,自己撞上
我注意的那些。

这些数字本身已经雄辩,并且出人意料地印证了去年就已强加于我的那个印象:对我的著作、对我
这个人的\emph{葬礼},不是一个人的事业,也不是一个严格有限的集团(比如我离去前的学生,或我
那些“上同调学生”)的事业,而确实是一项集体事业——其水平,是“教团全体”的水平;或者至少,
是数学建制中曾经\dmark{}\origpage{1113}\textbf{见证}并参与了我1955至1970年间作为几何学家的著作之兴起与繁盛的那一
部分的水平。我1970年的离去,至少在这一片数学里,成了一个信号:对“Grothendieck式”
数学,一场即刻而严厉的\textbf{排斥反应}——人们把它感受为“女性气质的数学”的象征与化身
\fnote{850}{(*)}{关于对某种对待数学的进路风格的这类排斥反应,见“肌肉与肚肠(阳葬阴(一))”、“阴的葬礼(阳葬阴(四))”、“天赐的机缘——或封神”、“拒认(一)——或召回”、“拒认(二)——或变形”诸注(n\textsuperscript{o}s~106、124、151、152、153)。在“上涨的海……”、“九个月与五分钟”、“箭与浪”、“兄弟与配偶——或双重签名”、“阴,作为仆人——与新主人”、“阴,这仆人——或慷慨”诸注(n\textsuperscript{o}s~122、123、130、134、135、136)里,我试着勾勒“女性气质的数学”的某些显著特征,与互补的“男性”特征并列。}:在那
里,视野永远先于并启发着技术的层面;困难永远被化解,而不是被砍掉;与事物在表面参差之下的
深层统一的持续接触,让人每时每刻都能从偶然与次要的无定形团块里,认出什么是要紧的。与此同时,我的离去也是一个信号:一切概念性的工作骤然停摆——或者更确切地说,一切此类工作被
\textbf{褫夺公权},以“深化”为名,骤然打上嘲弄的烙印。

就这样,砍掉数学创造工作的一个本质“侧面”——“阴”的或“阴柔”的侧面——人们达到的是一种
惊人的“扁平化”、“削平”、“干涸”\fnote{851}{(*)}{关于这种“削平”的初步确认,见注“无用的细节”部分(c)“什么都不像的东西——或干涸”(n\textsuperscript{o}~171(v))。}。这件事(在我看来)是通过一个骤然而严厉的急转弯完成的,几乎一
夜之间。这件事如此奇怪、如此闻所未闻,以至于显得难以置信。我花了一年多对\emph{葬礼}的高强
度思考,才终于把握住发生了什么,才认了这份明摆着的事实。我不知道,近年或近几十年,或任何
别的时代,在科学的一支里、或任何别的动用我们创造之力的人类活动里,可曾有过可相比拟的转
向。

可我回到我的相册上来。除了方才点到的诸位,把另一些人对\emph{葬礼}的参与——在我看来毫无疑
义的——的名字收进来,我想是有用的。况且,我并不确信他们中哪一位对我怀有恶意;他们之中肯
定不止一位,\dmark{}\origpage{1114}甚至对我怀着同情乃至好感(回应着我自己心里相似的情感)。他们中也许
没有一个人,会不真心错愕地听说:竟然有一场“葬礼”在对我这个人、对我的著作举行;更不用说
听说自己据说还以这样那样的方式参与其中。在这里指名道姓,本身就会产生那个(对我而言受欢迎
的)效果:把这事告知他们;并且(如果他们自己关心)由此给我们之间一个说明的机会。当然,关于
我所感知到的——无论对错——他们对我的葬礼的直接参与、或经由“共殓者”的参与,我随时听候
诸位垂询,提供一切细节。我绝无意质疑他们中任何一位的善意与职业操守
\fnote{852}{(*)}{(6月16日)随着新近传到我这里的信息,这份善意推定,在Borel的情形里化为乌有了。照他去年与Zoghman~Mebkhout的通信——事由是Borel在苏黎世主持的一个 D-模理论讨论班——我早已知道:Mebkhout曾告知他,自己是那门理论(所谓“Riemann--Hilbert”)中心的范畴等价的作者,并给出了确切的参考文献,还寄去了他的全部著作,Borel本可轻易从中确证事实。这并没有妨碍Borel以照例的居高临下(甚至失礼)对待他。在刚于奥伯沃尔法赫结束的一场同题讨论会上(Algebraic~theory~of~Systems~of~partial~differential~equations,~Oberwolfach,~1985年6月9--15日),Borel做了头三个引言报告(题为“代数 D-模理论”),为“好上帝定理”铺路;\textbf{Mebkhout的名字,在这些报告里一次也没有被提起}——在随后的报告里也没有(Brylinski的报告里,除了一句顺带的“喊停式引用”之外)。照Mebkhout刚寄给我的会议纪要,这场讨论会——Borel在其中指挥乐队(顶替没有出席盛会的Deligne)——是四年前那场反常层会议的名副其实的\textbf{重演}。“黑手党”几乎全员到场:Verdier、Brylinski、Laumon、Malgrange,甚至(这一回)Kashiwara(此人在苏黎世那场讨论班里,早已扮演头等角色,尽管Mebkhout已就其为人向Borel作过周详的通报)。不用说(同苏黎世一样),没人觉得有必要请Mebkhout做个报告;而且(除了这位Mebkhout偶尔的插话——它们都掉进了冰窟)先祖的名字一次也没被提起(除了那个不合时宜的“Grothendieck群”里到底有他的在场)。双对偶定理在那里依旧叫“Verdier对偶”,连Borel的报告里也是。Mebkhout早在去年就再三提请注意:这份双对偶,是从 SGA~5 的第一讲抄来的——看来Borel对某种风格、对某位缺席者,发展出了一种过敏症,禁止他考虑这一类参考文献……他自己在《交叉同调》(Intersection~Homology,Birkhäuser,1984)里,也成了同一骗局的参与者——而那本书,是在Mebkhout向他点破Verdier的把戏\textbf{之后}出版的。
对Borel,我曾把善意推定保持到极限:五十年代我们同在Bourbaki,一起共事,我对他知之甚深。在真正可称为我在数学世界的“出身圈子”的成员里,他是第一位,我不得不在今日、以无可怀疑的方式,确认其直接参与——而且是在“诈骗”的层次上——葬礼。},其中不止一位,我甚至可以补一句:
他们的全然的善意与诚实,于我是不容置疑的。\dmark{}\origpage{1115}与其愚蠢地按字母序开一张清单(那种
事计算机比我强),我宁可按大致的时间顺序,列出为我的葬仪齐声合唱的信众们的名字:不是按他
们在葬礼仪式上露面的时刻(那多半我并不知道),而是按我清楚得知其参与的时刻。另一方面,我
要把我全体学生单独列出\fnote{853}{(*)}{我在这里说“我的学生”,指的是跟我做到博士论文层级、并且(除Deligne外)跟我做了博士论文的那些人。共十四位(其中两位在我“离去之后”),见“耶稣与十二使徒”一注(n\textsuperscript{o}~19)的巡视。}。除了Hoang~Xuan~Sinh夫人——她在越南工作,分明显然太远,
帮不上我的葬礼——我的学生之中,没有一位,不曾以这样那样的方式参与其中。此事我已在“沉默”
那条注(n\textsuperscript{o}~84)与“棺材1——或知恩的D-模”(n\textsuperscript{o}~93)的开头说明过;这里不是重谈的地方。对
于我的每一个学生,我认为最值得做的,是一份就发生了什么而作的深入的说明。

“为我的葬仪合唱”,分布在极不相同的调门上。我辨出了四种主要的调门;它们为我说成了一场第一
流的、气派十足的多声部葬礼!有对一切带着Grothendieck气味的数学之发展的“沉默而有效”的\textbf{抵制};
有\textbf{失礼}与不体贴——我\dmark{}\origpage{1116}离去之前,在数学世界里我从未遇到过的;在一两个极端的情形,它们采取
了几乎不加掩饰的\textbf{嘲弄}的形状。有\textbf{刻意无视}或贬低我的思想与观点在其个人著作中的影响、
或在当代数学某一部分中的影响——而在那些情形,这影响分明而关键;或把无可置疑地属于我的
结果与观念,安到一个第三者头上。最后,还有那种(“\textbf{鸵鸟}”)姿态:碰巧面对一场刺眼的骗局
的人,把头埋进沙子,装作什么也没看见、什么也没感觉到。

不用说,信众的唱诗班里,不止一位在同时吊着好几个调门的嗓子。

话说到这里,终于可以兑现那份答应过的清单\fnote{854}{(*)}{我没有把那八位“非上同调”学生的名字收进这份名单;连同上面已经巡视过的上同调学生的名字,它们可在已引之注(n\textsuperscript{o}~19)里找到。
同样,在我的“家庭相册”里,也理应列出我的同事或旧友之中那些我确知为“不予埋葬者”的名字——他们以绝无含糊的同情与敬重之表示证明了这一点。首先,关于我1983年所做的那份工作《追寻栈》(我还打算回头谈这份工作),我收到过热忱的鼓励,来自 J.~Benabou、N.J.~Baues、A.~Jouanolou,尤其来自Ronnie~Brown与Tim~Porter——他们(不止以一两种方式)在我做那份工作的整个期间给了我切实的帮助。
确实,这几位同事所属的环境,与我惯常自认同属的那个环境相当不同;而我那场气派十足的\emph{葬礼},也正是自然而然地发生在后一个环境里。作为属于或接近这一环境、并且最近(在过去一两年间)给了我同一方向之表示的数学家,我乐于在这里点出 B.~Lawvere、J.~Murre、D.~Mumford、I.M.~Gelfand,以及(并非最末!) J.P.~Serre。这最后一位,独享同时出现在两份“名单”上的殊荣——“埋葬者”的名单,与忠实朋友的名单!}了——它来充实我们的家庭相册:Eckmann、Dold、Campo(N.~A.~Campo)、Mazur、Poenaru、Epstein、Cartier、Quillen、Kuiper、MacPherson(R.~D.~Mac~Pherson)、Heisuke~Hironaka、Hirzebruch、Tits、陈省身(S.~S.~Chern)、Artin、Langlands、Rota、Goulaouic、Fulton、Borel、Tate、Serre。

\subsubsection*{c.\ 众人之上的那一位——或首肯}
\tocsubsub{c.\ 众人之上的那一位——或首肯}

(6月11日)昨夜敲出那份名单时,我觉得自己有点傻:那里老老实实排开的每一个名字,单是它本
身,就唤起一整团丰富的联想,而在这里全无踪影。可这里绝不打算一一停在那些名字上、停在它们
所唤起的东西上——那需要再写一卷,而我正急着结束\dmark{}\origpage{1117}这一卷!我向当事的人们致歉:如此草率地把
他们“贴”进一幅不太能启人深思的“(出席我葬礼的)合影”里。确实,他们中的大多数,已在《收获与
播种》的这里那里以这样那样的名分出现过,即便未必是作为我的葬仪的参与者。其中有四位,属于
我Bourbaki小组的朋友圈子;三十年乃至更久以前,我便与他们因工作而紧密相系(其中两位更有友
谊的纽带)。在这份石碑式的名单里,另有九位,我自觉以温暖的友谊之情与他们相连;那份情谊,到
我写这几行时还没有熄灭。可是,逝去的岁月里,不止一次,当我又面对这些昔日朋友中的某一位、或
我昔日学生中的某一位,一种奇怪的印象攫住了我:我心头那份依然完好的、向对方奔涌的同情之
流,所奔向的那个人,已经不在了——或者说,与那个人的联系已经断了,也许无可挽回地断了;另一
个人顶替了我从前认识的那个人——那曾是个充满热烈而颤动的生命的人——而新来的这位,似乎
已把他的一切痕迹抹得干干净净。那像是一场\textbf{干涸},一场枯竭;在一度有敏感而活的血肉之处,
生出了一层又硬又不怕水的壳……

在这本我才刚刚半开的家庭相册合上之前,我还想在匆匆插入相册的众人之中,单挑一位,稍稍多
停留片刻。他在这本相册里来得最晚。比之于我最终收进来的任何别人,在我这里都曾有(当然是不
自觉地)严肃的抗拒:不肯抛开一些关于我们关系的陈旧的现成图像,不肯认一份谦卑的明摆着的事
实。我说的是Jean-Pierre~Serre。

在《收获与播种》里,我不止一次有机会谈到Serre,通常指名道姓\fnote{855}{(*)}{在《愚妄与更新》里,有两三次我忍住没有点Serre的名;在那些时候,对我多少持批评态度的人,通常并不需要点名。《收获与播种》中我谈到Serre、谈到我们之间关系的最周详的段落,在“九个月与五分钟”、“兄弟与配偶——或双重签名”、“无用的细节”三条注里(n\textsuperscript{o}s~123、134、171(v))。}。这里那里说的那一点,想必已足以让人感到:他
在我的数学家过去里,扮演过一个无人可以相比的角色。这件事,在写\dmark{}\origpage{1118}《收获与播种》之前,我自己
也从未在上面停留过;它是我一路写下去才发现的。二十年间——从五十年代初,直到我告
别数学舞台的那一刻——他为我扮演了“得天独厚的对话者”的角色\fnote{856}{(*)}{1965到1969年间——Serre与我的关系依旧紧密——充当得天独厚的对话者的,更多是Deligne。缘由,想必在于性情上的强烈投合,尤其在于Deligne的一种开放态度(对我感到是我所当贡献者之精华的那种开放),而Serre常常缺这个。我下文还要回到这两种关系的截然不同的性质上——它们是我数学生涯中最紧密的两种关系。另见下一条脚注所引之注。};我的许多大的执念观念、大的投入,都是直接被Serre的思想(有时“貌似平淡”)激发
起来的。有某些时刻——我印象里尤其是五十年代下半叶,也许还有六十年代初——他与我之间有过
一种强烈的数学“共生”;我们的数学性情是互补的\fnote{857}{(*)}{关于这种互补性,以及Deligne与我之间的投合,见已引注“兄弟与配偶——或双重签名”(n\textsuperscript{o}~134)。};这种共生,每次都证明极其丰饶。Serre与我之间的关系,并不是“对称”的:比如,
Serre完全不像我这样,倾向于托付给一个或几个“得天独厚的对话者”,来了解他可能感兴趣、或他自
认为需要的东西。尽管如此(至少我揣想),我在他的数学家过去里,想必同样扮演过一个非同寻常
的角色;我不妨想象,我1970年那次不辞而别,在他的数学生活里,是一个断点(也许是某种平衡的断
点——在那平衡里,我代表着“阴”的一极),是一次骤然的转向,凭着一股突然出现的“空缺”。我不知
道……

不管怎样,Serre同我这个人、我的著作之间的这层紧密关系,数学界想必是知觉到的,哪怕它一直
在“不宣之列”。可以肯定:除Deligne之外,Serre被(有理由地)看作与我的著作“最亲近”的数学家。Deligne同我的著作、我的关系则全然不同——那是一种学生与“继承人”的关系。Deligne以我的思想、
我写下与未写下的著作为养分;而我的任何大的执念观念、大的投入,没有一样是由他唤起或激发
的。在某种意义上,他比Serre“更亲近”我:在他与我接触的那些年(1965--69)里,他身上没有Serre身
上的那种对我的著作、对我对待数学的进路的某些侧面的排斥反应;正是这一点,让他得以在仅仅\dmark{}\origpage{1119}
三四年之间(凭他出众的才能,也凭格外有利的种种机缘),把我此前数年诞生、发展的那套宏大的
统一视野,由内而外、整个地吸收进去。但他同我的关系,根子上是深深的暧昧——而且
他系统地玩弄着这层心照不宣的学生与继承人的关系:那对他是一种\textbf{权力}的手段;同时他拒认
这层关系,忙于把宗师、把宗师的视野一并埋葬……

Serre与我之间,没有这一类的暧昧——无论哪一方,任何时候,都不曾有过要从对方身上攫取“权
力”的念头、或把关系用作权力工具的念头。我甚至可以说,在我四十年代末加入的那个“Bourbaki
圈子”里,这类权力游戏并不存在;直到我1970年离去,我不曾目睹、更不曾(哪怕身不由己地)参演
过这类游戏\fnote{858}{(*)}{不过,我应当作一点保留,顾及某个完全在我不知不觉之中、在我某些学生之间、围绕我这个人我的著作进行的游戏。那游戏至少从1966年(SGA~5 讨论班结束之年)就开始了;第一个清晰可见的插曲,是Deligne1968年关于谱序列退化的文章(关于这一点,见“驱逐”一注,n\textsuperscript{o}~63)。我开始知悉这些把戏——它们确实是权力游戏——是去年,近二十年之后。诚然,积极的演员,并非当初接待我、我也自认同属的那个圈子的成员(在那个圈子里,即便凭着成熟给我的回顾,我依然看不出这类把戏)。他们组成了“接班人”。同样确实的是:我在这些接班人身上察觉的、相对母体圈子的品质败坏,与那个初始圈子的几乎每一位成员身上(几乎无一例外)悄悄发生的类似败坏,必然紧密相联。关于这一点,见“Bourbaki,或我的大机遇——及其反面”与“自深处”两节(n\textsuperscript{o}s~22、23)。}。另一种说法(关于Serre与我的关系,或关于我在Bourbaki圈子里
观察到的关系)也许是:任何时候,我在那里都没有察觉到丝毫\textbf{对抗}的成分——哪一方都没有
\fnote{859}{(*)}{不过,七十年代初“生存与生活”的插曲,应当算作例外。那一插曲让事情大白于天下:在许多我以为重要(今天依然如此)的问题上,我自己的伦理与意识形态抉择,同我的数学建制朋友们的抉择——Serre也在内——几乎全体处于对跖点。正是它,骤然终结了我同那个“建制”的认同感——我那时倾向于把它混同于一个理想(且田园诗般)的“数学共同体”。(关于这一点,见“‘数学共同体’:虚构与现实”一节,n\textsuperscript{o}~10.)这场意外的揭示,以及数月之内由此而来的“改换阵营”,当时使我采取了对某些旧友的敌对姿态;我从此把他们归入“反动派”等等。后来,我从这些斩钉截铁而肤浅的分类里退了出来。话虽如此,经过一个并不奇怪的翻转,Serre属于那样一些人:有一阵子,我把他们感知为“对手”,甚至“讨厌鬼”。我高兴地确认:这段插曲没有在他身上留下怨恨或敌意的痕迹——在我身上,也丝毫没有,这还用说么!}。偶有摩擦\dmark{}\origpage{1120},这确定无疑;前面谈到过,或许还要回头谈;但那是全然另一回事。Serre与我的关系之力
量,在我印象里,恰恰来自我们对同一位情人——数学——的共同激情,其间不掺任何唯我的“寄生”
成分:在其中,对方不是手段、不是工具、不是靶子。所以,最近当我同Serre恢复中断了十来年的通信,
我从他给我的两三封信的字里行间,重新找到完好无损的友谊与体贴的迹象,仿佛我们昨天才刚分
手。

再说,即便十多年没有机会通信,Serre偶一传来的音讯,全都朝着同一个方向:不变的友谊——而不
是葬礼的调门,后者是我许多旧友的情形。所以,直到最近的这几个星期,我都未曾想过:Serre会在
我的葬仪上扮演角色。凡从他那里传到我这里的一切、凡我知道的他的一切,似乎都朝着相反的方
向。而且确定无疑:单是他在数学舞台上的在场,就曾给\emph{葬礼}划定过某些界限(不得不承认,那
界限质朴之至……)。翻阅 J.S.~Milne那本《平展上同调》\fnote{860}{(*)}{普林斯顿大学出版社,普林斯顿,新泽西。正是这位 J.S.~Milne,两年之后,参与“难忘文集”《数学讲义》第900种的骗局(见“……以及掘墓”一注,n\textsuperscript{o}~168~(iii))。}(1980年出版,故而在那场难以置信的
“SGA~4½--SGA~5 运算”\textbf{之后}),我惊讶地看到:Milne“放心地”、几乎逐字逐句地沿用了Serre在
某次Bourbaki讨论班(1974年2月,n\textsuperscript{o}~446)上关于平展上同调之父的说法——理论“由Grothendieck发展,\dmark{}\origpage{1121}得 M.~Artin之助”\fnote{861}{(*)}{两年前,1978年赫尔辛基国际数学家大会上, N.~Katz(还是那个Katz)颂扬新科菲尔兹奖得主Pierre~Deligne的演说里,平展上同调理论被说成“由 M.~Artin与 A.~Grothendieck发展,方向为Grothendieck所预定”——可见字母序有时候还真管用……Milne在诸说之中选择沿袭Serre的说法而不是Katz,在我看来,是其善意的又一个征兆。}。有种种迹象表明,Milne只在 SGA~4 与 SGA~5 里零星翻过
\fnote{862}{(*)}{尤其令我惊讶的是:Milne(Mebkhout也一样——他可是我的著作的一位细心读者!)没有察觉 SGA~5 里有一条显式的Lefschetz公式,关于代数曲线上一般的上同调对应;那公式,被魔术师-骗子二人组Deligne与吕西耶妙手变走——说这是艺术家的活儿,再贴切不过!关于这一点,见“变戏法者——或飞走的公式”与“恭贺——或新风格”两条子注(n\textsuperscript{o}s~169$_8$、169$_9$)。};他一面沿袭Serre(在同一篇Bourbaki讲义里对 SGA~4 与 SGA~5 满不在乎地发议论),一面又沿袭Deligne(在其笔名“SGA~4½”的锯倒之作里,对这些讨论班肆无忌惮地拆台),为的是在他的引言里,把原
始文本 SGA~4 与 SGA~5 呈现为“难以进入”\fnote{863}{(*)}{至于 SGA~5 的出版版——托“编者”(原文如此)Illusie的“福”,它只是原讨论班一座毁容的废墟——Milne觉得它“难以进入”,情有可原。好撒玛利亚人Illusie竭尽全力(顺着好撒玛利亚人Deligne的心意),把它做成了一本难以消化的“技术性枝蔓”大杂烩……}——而这,恰恰是他的书(继Deligne三年前那本——说实话,边上
有点单薄——之后)理应来补救的情形;说得明白些:免得用户去费那无用而烦人的工夫,去读原始文
本。最高权威们的意见(这里首先是Serre,继之以Deligne,外加一位待在软壁棺材里一动不动、一声不
吭的死者……)——意见,一个Milne如每一个平常人一样闭着眼睛跟从(在殡葬的语境下,甚至抢着跟从……)——斩钉截铁地排除:这些文本会是别的什么,只会是“无用的细节”(或“废话矿
渣”……),而绝非一门新的“广义拓扑”(topos 版)的基础——这门新拓扑曾使(除其他理论外)Milne
书中所论理论的绽放成为可能;长远看,人们没法省掉这门新拓扑,正如没法省掉寻常的广义拓扑
——那种拓扑,Milne、Deligne、Serre(同我一样)有福气在学堂板凳上学过,所以他们顺从地(当作
理所当然)承认:那门拓扑的戏,值得那根蜡烛……

\pmark\origpage{1122}我想是去年,我才第一次匆匆看了Serre那篇Bourbaki讲义;我最近在“无用的细节”(n\textsuperscript{o}~171(v))
部分(a)“一摞摞千页的大书……”里谈到过它。Serre讥讽 SGA~4 共1583页的那一段,当时给我留下的
印象如此之淡,以至于一两个月前,为写《四种运算》重新拿起同一篇讲义时,我竟把它忘了个干
净。应当说,Serre对我那些出名的“千页大书”的这种拉开距离的态度,我早已深知——早在 SGA~4
系列出现之前——所以它毫无让我意外之处。我记得,这种“发自本能的排斥”第一次在Serre身上发
作、冲着我对待数学的某种进路,是围绕凝聚对偶理论——我在五十年代下半叶发展的。那才是潜
在的“千页大书”(至少),尤其若算上随附的一整套新的上同调代数——导出范畴版;但“大书”也罢,
潜在也罢,明摆着的一点是:Serre听都不愿意听,正如韦尔受不了看人白纸黑字写一个上同调群、
听人说出“拓扑向量空间”这几个字一样。

\pmark 然而这一回\fnote{864}{(*)}{其实,这已经只是我第三次手里捧起这份文本,它才“终于让我心头一动”。},当我重读Serre1974年的这份文字——底色则是一整年关于某场\emph{葬礼}的思考(那
葬礼,1974年已经“大行其道”四年了……)——那一段终于让我心头一动。它在我心中不动声色地起作
用,一天天、一星期星期。我意识到:Serre的那种态度——我最终习惯了它、它在我离去之前“并不
牵扯后果”——竟像是对那场葬礼的\textbf{一盏绿灯}。这一方向的头一件事,以明摆着的力量向我显形:
Serre的措辞本身(但“添上了恶意与无耻”),仅仅三年之后,就被一位Deligne急不可待地接了过去(或
更确切地说,怀着隐秘的快意),当作他那场难忘“操演”的“背景噪声”。

我在已引的那道注——5月4日——里第一次这样表态;这场思考在同注部分(c)(5月27日)“什么
都不像的东西——或干涸”里得到深化。那里,也是我第一次开启关于Serre与我之间关系的思考——\dmark{}\origpage{1123}
以\emph{葬礼}所提供的特殊照明\fnote{865}{(*)}{在前面一条脚注(打字稿第1117页注(*))里,我还点过另外两道注,其中我谈到过Serre与我的关系,但照明相当不同——那是“葬礼之前”的照明。}。写那几页时,我心中想必已经有一种弥漫的知觉:Serre在\emph{葬
礼}里扮演的关键角色。此后两星期,一场对整整一批事实与印象的整合与消化,想必在无声中进行;
抗拒对事物直接而细腻知觉的惯性之力,我相信,已不经战斗、不费力气地消解了。时机,在我看来
已经成熟,可以把知觉到的东西表述出来,尽我所能。

人们也许会想:Serre由来已久的、对我著作的某些侧面与部分拉开距离的倾向,不过起了一种阴差
阳错的偶然作用——可惜,恰好促成了一场同样阴差阳错的葬礼。然而那是一种肤浅的看法,碰不到
事情的深处。径直说到问题的核心吧:鉴于Serre同我这个人、我的著作之间独一无二的关系,鉴于他
对其同辈、以及其后各辈数学家的非凡影响,有一件事在我看来已经清楚:\textbf{若没有他对我葬礼的
心照不宣的首肯,这场葬礼原是举行不起来的}。

除了一位分明缺席的“死者”,这场\emph{葬礼}还有\textbf{两位主角}:他们的作为与不作为互相衔接、互
相补足,似乎没有丝毫摩擦与涂痕(但在我这里,并不打算谈什么串通——两位主角的调门实在太不
相同):他们是Pierre~Deligne,与Jean-Pierre~Serre。

关于前者,在这场关于\emph{葬礼}的长考的一开始,就已有长篇的交代:他代表着葬礼图景的“前景”,
作为葬仪大主祭,同时是那些由他倡导的运算的隐秘继承人与主要“受益人”(而且,早在死者的象征
性“死亡”之前就如此……)。Serre——这里第一次作为葬礼仪式的前排人物谈到他——则代表着图景
的“第三平面”,由“信众教团”构成的那一平面。

\pmark\origpage{1124}早在去年——或者更确切地说,早在我以它最粗粝、最悖谬的形状(以及这个名字)发现\emph{葬
礼}之前——我就清楚地知道:在一个我从不曾知道自己有仇敌的世界里,以这样的热忱来埋葬我的,
首先是\textbf{我昔日的}\textbf{朋友们};他们中的某些人,却并没有因此停止(哪怕只从唇边……)把自己算作
我的朋友——如今,我同样看清:在这些朋友之中,那些同时(而且尤其)是我的学生的人之外\fnote{866}{(*)}{在《收获与播种》的整个思考里,越来越清楚地显出:做过某个人的学生(比如做过我的学生)这个单纯的事实,就给一层关系打下烙印,赋予它一种特殊的质地,使它接近与父亲或母亲的关系。},真正
构成那场仪式的\textbf{台柱}的——作为教团的代表、作为全体信众首肯的担保人——正是他们中间
那位,在我们共同激情的层面上,曾与我最亲近的人。

Serre首肯的最耀眼的标记,当然不在我所熟知的那句随口的俏皮话里——那句话差点从我眼皮底
下溜走(虽说并没有逃过所有人的眼睛……)。真正令我驻足之后惊骇欲绝的标记,是\textbf{他所乐于
保持的蒙昧}:对那场就在他眼皮底下进行的葬礼——这样说毫不夸张——蒙昧到最后一刻
\fnote{867}{(*)}{可以说,在他那篇已引的1974年Bourbaki讲义里——他陈述Deligne对Weil猜想最后一角的证明——Serre的鼻子正埋在葬礼的正当中,却天真得连想都没想过记下一笔。我隐约感到他的不安:他发现自己面对一种表面荒谬的处境:我作报告(也在Bourbaki讨论班)给出某条 $\ell$-进上同调的 L 函数公式之证明的大纲,十年之后,我在其中承认(当作已知)的那条关键的“不动点公式”,居然仍未在文献中得到证明。Serre于是选择用一阵火气来\textbf{排遣}这份不安:讥讽 SGA~4 那出名的“1583页”(言下之意:而且\textbf{连}人们需要的那条公式\textbf{都没}提供)。这是一条省事的路:把一件不愉快的现实一笔带过(x)。然而他心里非常清楚(也许他乐于忘记……):在 SGA~5 讨论班里,我曾原原本本证明过一条不动点公式,其范围远远超出对 Frobenius 对应的那一条;他也知道,我的讲稿的成文,拖在那些所谓的义务“编者”手里,已经八年了。他即便乐于忘掉 SGA~5 的主题(“L 函数与 $\ell$-进上同调”——这标题好歹还是说了它该说的)及其内容,可他认识我二十多年,看我做数学看了二十多年,深知我把事情做半截不是我的习惯——恰恰相反(而且我做事情“不半截”到这样的程度,常令他不胜其烦,乃至忍无可忍……)。这本可以帮他刷新对 SGA~5 讨论班上所发生之事的记忆——他的脚,反正常踏进那扇门,足以大致知道我在那里做什么、进展如何。
显然,他既不想刷新记忆,也不想提问题。这只是众多例子中的一例:我的朋友宁可闭眼、塞住鼻子,也不肯去知悉一种他若接受就必然深深牵连他自己的真实。
(x)(6月22日)这几行写下之后,我有机会确认:这一类“不愉快的真实”,如今竟受到热烈的欢迎,如同天降的横财!关于这一点,见“家庭相册”的 d.、e. 两部分。}——那是一场对一部著作的葬礼;\dmark{}\origpage{1125}著作自其源头起就与他相系,比世上任何人都系得紧。而读《收获与播种》(假如他肯读这“一大摞”上千页的书的话……)会不会终于促使他动一动自己的鼻子(十五年来,馊味早就很大了……),以及其余的一切——对我而言,这是一个彻底的谜。可我深知:对他来说,同对葬仪的任何其他参与者一样(甚至更多),接纳我的讯息、动用他健全的官能,也意味着接受让他自己被深深牵连进去。

Serre的角色——率领那支前来观礼、为我葬仪合唱的信众教团——在我看来既典型,又例外。说它
例外,一是其极端性:他是与我最亲近的人,比教团任何成员都近;二是其非凡的身量(器局)
\fnote{868}{(*)}{还有第三种情况,赋予Serre在葬礼中的角色这种“例外”或“极端”的性质。他属于在我初次接触数学家世界时接待我的那群“善意的长者”。(关于这个群体,我有生以来第一次在“受欢迎的异乡人”(n\textsuperscript{o}~9)一节、继而在《收获与播种》引论(I~5,“一笔情愿的债”)里谈到过。)除了我们之间的友谊与好感之外,这也许是最主要的缘由,使得我花了一年多,才认了那份明摆着的事实,确认了Serre在我的数学葬礼中所扮演的关键角色。}:这
身量(器局)排除了深层动机中“补偿性的对抗”这类惯常的“寄生”成分
\fnote{869}{(*)}{这种“(表面上的)无因对抗”,我在这里那里已影射过两三次,尤其在(4月3日)“信使(2)”(n\textsuperscript{o}~182)一注(见下文)里。对我葬仪的绝大多数参与者而言——也许甚至是除Serre之外的全体——这种“原型式”的对抗肯定在起作用,这一点,我毫不怀疑。这股力量,在我看来,有别于那股通过“被自己否认的女人”之压抑(或“埋葬”)过程表达的力量。但这两种力量又紧密相联;在葬礼中它们结成一伙,以某种合金出现,常常难以分开。不过,我相信我已在其中辨认出葬礼中起作用的\textbf{两大力量}。可眼下,我很难说其中哪一方占先,是哪一方。我倾向于认为,是我先辨认出的那一方,即:对自己存在中阴柔侧面的压抑之力。
Serre的情形之所以在我眼前显得“典型”(同时又例外),想必是因为,在场两股力量中,正是后一种(我倾向于视为根本的那一种)在其中显出其全部力量,而丝毫没有另一方的痕迹(我称之为“寄生”的那一种——其义为:它会遮蔽对我所相信知觉到的\textbf{要害之物}的清晰把握)。我揣想(只要对原始事实与知觉的整合消化工作继续下去),未来几个月会带来更细腻的理解:在葬礼中、也在我以这样那样名分参与的其他冲突处境中,两股力量各占几分。}。如我方才强调的,我在\dmark{}\origpage{1126}Serre同我这个人、
我的著作的关系里,察觉不出丝毫对抗的痕迹;很清楚,在他的首肯之下起作用的深层力量里,也没
有对抗的痕迹。据我所知,除了那句出名的俏皮话,这份首肯纯然被动地表达:不过是一些\textbf{遗漏}
罢了。可是,这道给一场规模浩大的葬礼的心照不宣的“绿灯”——伴随的那些行动,有时庞大到仿佛
既定义了常识、又定义了体面——如今在我看来,恰是Deligne对同一场葬礼那积极而热切的参与的
不可或缺的\textbf{对应物}、某种意义上的“底片”\fnote{870}{(*)}{这里,Serre与Deligne在葬礼中的角色分配,有一处相当显著的\textbf{颠倒}:Serre的角色几乎全然被动,Deligne的角色则极度主动(虽说这位“操盘手”的角色,为事势所需、也合我朋友皮埃尔特有的风格,始终隐而不显)。然而,事实上,正是Serre其人,具有强烈得鲜明的“阳刚”主导;而Deligne其人,具有同样鲜明的“阴”(或“阴柔”)主导;并且(对两者皆然)无论在唯我机制的层面、“自我”及其条件作用的层面(亦即“\textbf{老板}”的层面),还是在发现之冲动的层面——那本源的、就其内在本性而言逸出条件作用的层面(即“\textbf{孩子}”的层面)——皆然。在Serre与Deligne这对截然相反的极端性情之间——葬礼的两大“台柱”——死者本人,则代表一种中间项:在“老板”一侧,阳刚主导强烈;在“工匠”(或“孩子”)一侧,阴柔主导同样强烈。(这一“基调”的分布,出现在注“兄弟与配偶——或双重签名”(n\textsuperscript{o}~134)里。)
阴阳角色之间的“对调”之力与机制,本是\textbf{主要}的思考主题:它催生了“阴阳之钥”这场长考,并作为底纹贯穿始终。它从《钥》的第一道注“肌肉与肚肠(阳葬阴(1))”(n\textsuperscript{o}~106)起就隐含在场,并在后续十一道注(n\textsuperscript{o}s~124、127、132、133、138、140、145、148、151、153、154)中或多或少走到注意的前台。而在这里,我意外地撞上了一个多少类似的“对调”处境——其动因,是葬礼中起作用的深层力量的内在逻辑。
近来,另一处似乎悖谬的阴阳角色“对调”,也打动了我——在这场不乏表面悖论的葬礼里!这一回,说的是“过早的死者”一方,与葬礼全体参与者一方,各自的角色。在集体无意识的意图层面上,这场对一个死者(按理,他只应安于与其身份相称的完全被动)的葬礼,首先是一场“女性气质的数学”的葬礼——一种带着强烈“阴柔”意涵的数学风格与进路的葬礼;而埋葬的教团,按理体现着“纯粹而坚硬”的阳刚价值,把那软塌塌的阴柔酥软之物交付给应得的蔑视。(比如见“阴的葬礼(阳葬阴(四))”与“天赐的机缘——或封神”两注,n\textsuperscript{o}s~124、151.)然而,处境的内在逻辑,却逼着这些“纯粹而坚硬”的参与者,在其中演一出典型的“阴”的或“阴柔”的戏:一场“天鹅绒爪子”的戏——半明半暗、沉默、遗漏、不动声色安插的暗示,或不断地暗示一事而佯言其反——总而言之,“喊停!”的那一套;我朋友皮埃尔是其中的至圣至师,而每一位埋葬者都不得不多少把它据为己有,事势使然。(关于这套风格,见“喊停!”一注,尤其“天鹅绒的脚掌——或微笑”与“对调(四)——或夫妻的马戏”两注,n\textsuperscript{o}s~77、137、138.)相反,是那位“死者”——那丰沛的阴柔酥软之物的化身——在人们最不提防的时刻,从舒适的棺材里坐起来,骤然接过一个他熟悉的“男子汉”角色:摊牌亮底,把好奇的鼻子和不逊的舌头、连同手电筒,伸进那些最精微暧昧的昏暗角落,粗声粗气地逐个点名,是猫说猫,是贼说贼——总而言之,一个真正的没规矩的家伙,一个不折不扣的搅局者,搅的正是那一团和气、细语呢喃的安稳局……}。而我朋友身上发生的这场变化,分量惊人。

\pmark\origpage{1127}我似乎确实深切地感到了在Serre身上起作用的那股力量。它所处的层次,比一种个人之间的对抗、或比谋求某种“好处”(按这个词的通常意义)更深。近来与他交换的那几封信,在这一点上很说明问题。说来,十五年来——自我离去以来的十五年里——我的朋友身上完成了一场\textbf{转变}
\fnote{871}{(*)}{“转变”(\textbf{transformation})这个词,立刻把我带到我朋友皮埃尔身上的“\textbf{变形}”(métamorphose)——那是去年10月他来访时,我第一次清楚知觉到的。(关于这一点,我在“拒认(二)——或变形”一注(n\textsuperscript{o}~153)里谈到。)“变形”一词更有力,对应着这一事实:在我朋友皮埃尔身上,发生了一次真正的翻转——从一种“阴”主导鲜明的本来性情,翻转到一副借来的、“彻头彻尾男子汉”的姿态。撇开这一点,我在两位朋友身上感到的转变,方向相同,动力相同:那股压抑“阴柔”特质的力量。},方向恰恰是那种对我对待数学的进路中的某些主导侧面的“发自本能的排斥”
\dmark{}\origpage{1128}。这些侧面,当年在Serre自己的进路里也曾存在,只是程度较轻——那是在他
作为数学家的过去里最丰饶的那些年:开放的、创造力澎湃的年月,尚未有一道\textbf{压抑}其创造个
性中这些侧面、压抑其内在“孩子”的工序立起来之前。这些,是创造力的“阴”的或“阴柔”的侧面与特
质。我在朋友身上以惊人的力度感到的转变,就是:从阴与阳两股创造之力和谐协作(而“阳”或“男
性”主导鲜明)的状态,转为一种“彻头彻尾男子汉”的失衡状态——“阴”的或“阴柔”的品质被毫不留
情地连根拔去。

说实话,正如两周前我已暗示过的(前引之注),这是一场演化的终点;五十年代,我已辨出它的最初
迹象,六十年代里它愈演愈烈。那时起,平衡就逐渐破裂,表现为视野的\textbf{窄化},表现为获准
上场的那份创造才能的清单的\textbf{缩水}。对我对待数学的某些主要侧面的排斥反应,以及渐渐地,
对一切真正构成我的著作的生命、深度与力量的东西的排斥——那不过是向外投射:是在他与我的
关系的层面上,一份远为深重的排斥的可触显形;被排斥的,是他自己存在的一个本质侧面,是他自
己的创造才能的一个本质侧面。

也许(如我方才暗示的),只要我还在跟前,同我的关系就对Serre的这一演化起着刹车的作用;在
他的生命里——五十年代,尤其六十年代——它代表着某种配重,由此成为一个相对平衡的因子。果真如
此,我的骤然离去,必定给那股压抑阴柔品质的力量放了行——那类力量,我并不陌生:它
也是作用于我自己生命的主导唯我之力之一;只是有一个显著的差别:在我这里,这份压抑之力只局
限于唯我机制的层面、我与他者的关系的层面,\dmark{}\origpage{1129}从不干涉我同数学夫人的恋情,也从不干涉(更一般
地)我在发现历险中的自发举步——无论数学的还是别的
\fnote{872}{(*)}{关于这股压抑之力在我自己生命中的角色,我在“超级爸爸(阳葬阴(二))”(n\textsuperscript{o}~108)一注里谈到过。我是在1976年开始察觉这股力量的;那一年,是我的心灵历险的一个关键转折。那转折,在“重逢(阴的苏醒(一))”与“接纳(阴的苏醒(二))”两条注(n\textsuperscript{o}s~109、110)里谈到。至于“阴柔”特质在我的数学工作中的主导地位(这些特质似乎在那里避了难,躲开了一切怀疑!),见“上涨的海……”一注,n\textsuperscript{o}~122。}。

回到\emph{葬礼}:眼下我做得最好的,莫过于引录11月10日那场思考的结尾几行——“阴的葬礼
(阳葬阴(四))”一注(n\textsuperscript{o}~124,打字稿第564页):
\begin{quote}
“……而这葬礼,骤然间以一种新的、意想不到的光呈现于我:连我这个人也变得次要了,我成
了那‘理应交付蔑视之物’的\textbf{象征}。这不再是某个人的葬礼,不再是一部著作的葬礼,甚至
不是一场不可容忍的异端的葬礼;而是‘数学中的阴柔’的葬礼——而或许更深一层,是在那众
多前来为悼词鼓掌的与礼者每一个人身上,\textbf{那活在他自己之内的、被自己拒认的女人的
葬礼}。”
\end{quote}
这最后一条直觉,是那一天、写下这最后两行的当口,以一道闪电突现的,像是一份意想不到的启
示,加在前几行所说的那份启示之上。此后几个星期,这道直觉一直作为底纹留在我的思考里,最终
在12月23日至26日连续的三道注里被重拾、深化:“天赐的机缘——或封神”、“拒认(一)——或召回”、
“拒认(二)——或变形”。

在这道直觉初次显现的那一天,以及在所引三道注的前两道——我在其中对它稍作探测——里,我
心里并没有一个具体的个案,除了(稍稍地)我的朋友皮埃尔(在第三道注里得到更周详的检视)。此
外我也深知:这一个案,对构成我\emph{葬礼}上出名的“第三平面”的信众教团之全体,绝无典型性。
所以,由于没有一个具体个案来作例证,我对那场突然瞥见的真实的把握,仍带着\dmark{}\origpage{1130}一层含混——被预感之物的那层含混:在某个层面上“知道”,却未曾充分而清晰地“看见”。我隐约
记得,我曾为这层含混稍感不安;曾有一阵冲动,想在我的朋友们——我知道他们参与了\emph{葬礼}
——找一个“有代表性”的,好把这份弥漫的认识“挂”上去,看它化身为一个可触的真实。

那时,Serre的念头,一次也没有掠过我的心头——他属于我旧友中那寥寥数人:我曾断然认定(至少
在意识的层面上),他至少不在参与我葬礼之列!而我那摸索着的思绪,当时(以及此前……)没有找
到那个在葬仪上须以一身体现“教团全体”的人,想必是因为,在我之内的某个地方,一清二楚:世上只
有\textbf{一个人}能演这个角色——而恰恰是这个人,我之内的一种沉重,让我以某种心照不宣而斩钉
截铁的禁忌,一开始就把他排除了……

如今,随着一场缓慢而幽暗的地下功课,那份沉重已经消散;现在,全光照亮了:恰恰是他——众人
之中恰是此人——那份寻找化身的直觉,在他身上以如此完美的方式落实,以至于人们几乎要相信:
正是他本人,让它在我之内浮现;正是他,从它显现的那一刻起,就赋予了它那种斩钉截铁、无可辩驳
的“已知之物”的力量\fnote{873}{(*)}{我甚至满心倾向于认为:这个“人们几乎要相信”,恰与事情的真实相合。这将再一次证明:我们的认识官能,远远超出那片可怜而微末的映像之域——我们只容许它们进入意识目光那片划得极窄的田地。(6月14日)前天的思考收尾的那道念头、那道骤然的直觉,也是在落笔之时以“闪电”方式显现的,没有明显的准备,也没有检视的打算。它带着某种“明摆着的力量”呈现在我面前。只是事后,我才记起:在紧挨着11月10日那段引文所出自之注的前一道注里,我恰好有机会相当周详地谈到Serre其人、谈到他与我的关系(而且,这还是《收获与播种》中的第一次)。}



\subsubsection*{d.\ 入殓——或自然的坡道}
\tocsubsub{d.\ 入殓——或自然的坡道}

\pmark\origpage{1131}(6月17日)现实,每一次都超出一切预感(无论多敏锐),甚至超出人们对它可能有的“知
识”;只有当我亲身与之摩挲——在路上的转弯处,而且往往全然出其不意——我才渐渐浸透它的滋味
与气味。即便这种接触似乎只是\textbf{印证}了已被预感或已被“知道”的东西,它却往往出人
意料,毫不客气地撼动一种几乎拔不掉的\textbf{怀疑}:对那确凿已知、已说、已写、已重述、又已重写的
东西的怀疑——尽管在某个层面上(那是巨大的沉重所在的层面),它依旧是一纸空文。我不止一次察
觉这份沉重\fnote{874}{(*)}{关于这份“沉重”、这份“不信自己健全官能之见证”,另见“义务已尽——或真话时刻”一注(n\textsuperscript{o}~163)第782至784页,尤其第782页脚注(**)。},我的不耐为之火起——一股顽强的沉重,执意要把我拽回熟
悉的观念与图像的车辙里,或拽回那些多少得到普遍首肯的观念与图像里;而我分明“知道”(或者说,
我之内的\textbf{别的}什么人、什么东西知道……)这些安顿妥帖的观念与图像是一场障眼戏,往往是明
摆着的障眼戏,根本站不住……思想,哪怕怀着要知道底细(那件既被“知道”又被否弃之事的底细)的
强烈渴望,单凭思想本身,无力抹去这份深深锚在自我结构之中的沉重。唯有与真实直接接触的斩钉
截铁之力,才有能力有时撼动这份沉重,侵蚀它、挪动它分毫——若谈不到真正抹去的话。

昨天我给Serre打了电话。只为一件求实的小事:关于Tate的“刚性解析空间”笔记——近日谈到的
\fnote{875}{(*)}{见“黑手党”一注(n\textsuperscript{o}~171)部分(c),“记忆的失灵——或新编历史”。}。我模模糊糊记得,那份文本原有一篇短引,交代这项工作的来历;我印象里,这篇引言在1971年
《数学创造》做的那一版里“丢”了。其实,Serre向我确认:Tate的笔记里,本来就没有这样的引言。那
不过是一些日志式的笔记,Tate把他关于刚性解析的思量寄给Serre,几乎像写信一样,而且(当然)
绝无发表的念头。我记得是我安排它们经法国高等科学研究所散发的(作者署名之下,加了一个小
标题“私人笔记,经(未经)本人许可刊布”\dmark{}\origpage{1132}),可我忘了Serre是中间人。反正,1962年刚性
解析空间诞生之时,除Tate与我,最“在局中”的就是Serre了。是他,或许早一两年,向我讲解过所谓
“Tate”椭圆曲线的理论——取一个完备离散赋值环的分式域 $K$。我被那——我记忆中是一场显式公
式的洪水(据说还是“经典的”)——砸得有点发懵;它们从我头顶掠过,没有“挂住”。但有一幅惊人
的几何图像留了下来,想必是Serre一句评点所致:归根到底,Tate椭圆曲线(至少它的“点”),是在
乘法群 $K^*$ 里“商化”、被一个同构于 $\mathbb{Z}$ 的离散子群相除而得。这正是复情形的类比:先把
$\mathbb{C}$ 除以一个因子 $\mathbb{Z}$ 得 $\mathbb{C}^*$,再除以一个因子 $\mathbb{Z}$,这次得到的就是一条椭
圆曲线。在复情形,商化在复解析的领域里有确切的意义,而Riemann-Serre型(即 GAGA 型)的定理保
证:最终的商(一条紧复曲线)具有典范的\textbf{代数}曲线结构。而在Tate的情形,唉,在 $K$ 这个完备赋
值域上那套让人不惯的解析空间里,得到的商是一个\textbf{全不连通}的紧解析空间,根本没有可能从
中取出一条椭圆曲线。然而(Serre当时必是这样对我说的)一切却仿佛……总之,Tate硬是靠着 $K^*$
和它的离散子群,凭着一手显式公式,造出了一条名副其实的椭圆曲线。

我记得清清楚楚:无论Serre还是Tate,都不相信Tate那套计算性的构造,会真有一个“解释”——以
$K$ 上一种“解析簇”的新概念来解释
\fnote{876}{(*)}{(1985年9月)正如去年7月我与Serre的通信所显现的,这里有一处我记忆的变形(Serre也有一处)。Tate(8月4日、10月16日)与我(8月18日、10月1日、10月19日)的几封信——都是写给Serre的——使事件的影片得以重排。是Tate(既非Serre,也非我)最先有了直觉与确信:应当存在一种“解析簇的新概念”,以便简单地解释“Tate椭圆曲线”的形式主义——那是在1959年8月前后。在我这里,它并没有“立刻让人心头一动”(照我记得的样子);我对Tate这一建议的第一反应相当怀疑,然后才开始思考这个问题。想必不久之后我就被说服了——一旦我察觉:既有的观念(尤其形式概形)无法解释与Tate椭圆曲线相联系的现象。此后两年,我相信,独自思考新概念的界定原理的,只有我一人;Tate与Serre都毫无头绪。这种情形持续到1961年10月:我向Tate提供了一门理论的总施工图。这立刻触发他去发展必需的基础,以便把握那些局部碎片(这项工作,若没有一个清楚的想法——日后如何把这些碎片拼装成整体的对象——本不会有太大意义)。更详细的评说与相关书信的引用,见《数学反思》第3卷预定的“历史评注”。}。

至于我,它当年是立刻让我心头一动的(*),我压根不曾想过,可以不把Tate曲线“看”作商化的结果
——为一个尚待提炼的合宜“簇”概念商化的结果——那正是让我着迷的那类工作!很可能也是Serre,
尽管满腹疑虑,向我提过:有一些人,至少克拉斯纳,在完备超度量赋值域——因而全不连通的域——
上“做解析延拓”。这似乎为我那个(有点怪诞的)希望添了水:不管怎样,应当会有一种“好的”解析簇
概念,比已知的更聪明,并且(凭“连通”一类的性质)接近实或复的、甚至代数的\dmark{}\origpage{1133}解析簇。可
再一次,在那个三人小组里,真正相信的只有我一个——至少我当时有这个印象。

这事在我脑子里转了好几个月,也许一年——说不清了。这处境让我想起一桩旧日的困惑:在当
时可用的概念框架里(靠赋环空间,概形与形式概形之类),人们无力赋一个形式概形——取所论的
离散赋值环 $A$——的\textbf{一般纤维}以意义。很快清楚:这本质上是同一桩困惑——我为给Tate的构造
以几何意义而寻找的那类“簇”,应当正是那种能给这出名的、尚不存在的“一般纤维”以意义的簇。
我终于有了第三条主线(克拉斯纳的风声之外):那是1968年出现的“广义拓扑空间”的直觉——当
时还没有得到“site”或“topos”这样的名字,因为我还没有开始边写边琢磨的概念工作——它将使我
得以定义那出名的、(隐含地)进入Weil猜想的“$\ell$-进Weil上同调”。这提示我:正如Weil上同调
的情形,我寻找的那种新“结构类”,不应到那老一套的寻常“赋环空间”里去找,或许要到那些“广义
空间”里去找——配以合宜的环层。

这些零零星星的直觉,何时才强大到足以推我开一个括号、从手头的任务里(主要是 EGA 与
SGA)抽出时间,动手做一点边写边琢磨的雏形工作——我说不清了。我知道的是:这项工作,一如既往
地,是在孤独中完成的——只有我一人“看见”那里有什么;也因此,只有我一人当时有能力做那把它
带到光天化日之下的最初的工作。我记得,我这里想几个钟头、那里想几个钟头,甚至有时一整天,
像逃学一般(虽然“日常”的工作并不缺!)。终于有一天,我咬紧缰绳,要弄个水落石出,真刀真枪地
干上了——我至少连续花了好几天,也许一两个星期。最难的是超越那些根深蒂固的思维习惯:它们
不停地要把我拽回已知的常规——“寻常”解析空间的车辙(我想,如今叫“松弛的”——德文作
“welk”)。我不得不重来三四次——一发觉自己又回到老辙里(马回马厩了!)就再爬出来\dmark{}\origpage{1134}。但分明,这\textbf{不是}那匹老马能顶用的时候……


这项工作做完,我心里有了底:再加一份我那时没有兴致去做的技术性工作,我已经立起了一个
“刚性解析空间”的观念(这是我给它的名字;“刚性”一词,为的是表达连通一类的性质——贴近代数
簇,而与所谓“松弛的”解析簇的性质恰成对跖)——它至少足以回应当时我心中的两个愿望:以这些
空间解释Tate的构造,以及解释形式概形的一般纤维。


我当时没有想得更远——我最急着回到那些被我暂时撂下的任务。倘若我再多玩一玩,我很快就会
发现:连闭环 $r\le r\le R$ 这样简单的空间(它们也配得一种“刚性解析”结构)都逃出了我的构造。是
Tate——我当然让他知道了我的思量——做了必要的调整,好把它们包括进来。除了我做的大半的
概念工作,还有一份更技术性的活儿要做:对所用的“砌块”(它们扮演仿射概形的角色)要有好的掌
握。这活儿,恰以其特有的优雅与细致,在Tate1962年的笔记里完成了\fnote{877}{(*)}{为求公允,我想可以说:我的工作与Tate的工作,代表刚性解析空间理论之绽放的两个阶段,同等不可或缺。我的那份,在于最初的视野(Tate与Serre都不曾有)与一份以概念为主的劳作——它也绝非没有必须迎面处理的技术侧面。Tate的那份,以技术为主,但其中同样有一份概念工作。我的工作,基调是“阴”的、“阴柔”的(所以,除了我已不在场,它还成了人人轻蔑的对象);Tate的工作,基调是“阳”,合乎好品味与好规矩的法度。\par 倘若我不曾把自己的工作推进到足以不再有任何怀疑的地步——怀疑一种好的“刚性解析”观念的\textbf{存在},以及对一门理论的总施工图的清晰视野——那么这门概念,恐怕今天还不会问世。因为,虽说人们“早晚”必然会发现并发展这个概念——它当然不是什么“发明”——其需要,在此后二十三年里,却从未以足够紧迫的方式出现,逼人“跨出那一步”。显然,是我第一个(1966年)预见到了刚性解析理论在最初的两个动机之外的应用领域:晶体上同调的发展。\par 此外,除我预见的三个用途——包括把Tate理论推广到一般的阿贝尔概形——我没有听闻别的几何应用。看来,后来“在这一行里工作”的人们,只从中看到关起门来发展理论的材料(既然理论已存在,而共识把它列入“严肃研究课题”),却不把它嵌进更宏大的几何视野。这是一个鲜明的案例:数学思想的原子化、碎块化——它与一切奠基工作的失宠相连,与一切不能归结为某种解“竞赛题”的技术绝活的失宠相连。一个特别雄辩的征兆是:无人尝试发展一个更一般的刚性解析空间观念——使Tate发展的那种,与“概形”之对“域上的代数簇”相仿——好把“变动的”完备赋值域上的诸刚性解析几何联系起来(尤其特征变动的,并包括实、复与“超度量”诸情形)。这一缺席,是近十五年数学在一切奠基工作层面上惊人停滞的众多征兆之一(这层面的关键,在此显然)。\par 回到Tate与我:同样可能的是,倘若我最初的“突破”没有在Tate那里“让人心头一动”、没有触发他的“第二回合”,刚性解析空间也不会存在。这事我原会在周围这里那里谈起的;可诱人的题目(包括看似更“紧迫”的题目)从来不缺,恐怕没有人会留意它——尤其今天,肯定不会:如今,引入如此怪诞之物的念头本身,就会让人想起某位,还是不点名的好……}

\pmark\origpage{1135}话说回来,还要过一阵,我才认了这份明摆着的事实:我所用的砌块,边边上短了点。它们对激
发我的那两个初始问题,已然够用——何必再往远里找!我很难改口。最终是Tate说服了我,以他那种
既平和又细致的方式:毕竟,不止这两个例子;就算我此生看来没遇到过圆环,也不构成排除它
们的理由。而显然,用我那套砌块,没有办法把它们“收编”(除非用无穷多块——那又差不多落回“松
弛”的老辙)。

我做好了我那份工作——在我自己的角落里,这很自然,反正我是唯一相信的人——不过当然,一旦走到(暂时的)终点,我也没有放过机会,把这事告诉两位主要当事人\dmark{}\origpage{1136}(实际上也是仅有的两位知情人):Serre与Tate。Tate那边分明心头一动;而Serre,我想,在我告诉他我的进展时,想必也被说服了。对此我没有确切的记忆,但倘若事情偏偏不是这样,我肯定会记得的。

\pmark 所以,昨天我给Serre打电话时,在我看来天经地义:他和我一样清楚,在这种新簇观念的诞生中,我扮演了什么角色。我没料到会有机会影射这件事;可当我跟他谈起Tate的笔记时,倒是他提醒我:那些笔记原封未动地刊在《数学创造》上,而且Remmert与另外两位作者刚出了一本专讲那些出名的刚性解析簇的书。这本书,就是我最近在“黑手党”一注部分(c1)(“记忆的失灵——或新编历史”)里谈到过的那本——我在其中痛斥Remmert“记忆的失灵”(可Tate的笔记本可以替他唤回记忆),服务于一种在我看来明摆着的蓄意刁难。我顺带向Serre提了一句——在给他写的上一封信里,我已有机会影射过某场葬礼\fnote{878}{(*)}{这是对那封信的复信里的话——那是我收到的Serre的最后一封信——其中Serre引了Siegel关于当代数学“扁平化”(Verflachung)的评语;我对它作了评说,并在“无用的细节”(n\textsuperscript{o}~171(v))部分(c)“什么都不像的东西——或干涸”里继续讨论。如我在该注所说,Serre曾把Siegel的这份印象斥为“不公正”——可我感到,Siegel竟这样想,多少让他心绪不宁。而这同一个词,他(肯定不是故意地)又拿来,打发我对一场葬礼的影射。不用说,他压根没有想到要问我:\textbf{什么}使我说有葬礼(我在信里一个字没提,宁可等他来问)。显然,案子早已审理完毕……},而这毫无疑问是一幅相当触目的写照。

头一件相当疯的事,是Serre(天知道他当年可是在最前排的!)——唉,连他也不记得了,而且是压根儿不记得了:我在那些出名的刚性解析簇里曾经有过一份!我简直瞠目结舌!这真叫人发疯——当我向他影射我以为自己有过的那份微薄份额、从激发我的那两个例子说起时,Serre记得的恰恰相反:差不多是我压根儿不想知道这些新簇,说(照他讲)有了形式概形,就已经有了所需的一切\dmark{}\origpage{1137}!当下我简直不敢相信自己的耳朵\fnote{879}{(*)}{事后回想,我明白了朋友的记忆(边上略有些失灵)发生了怎样的变形。既然我曾把形式概形当作主要的、几乎是唯一的向导,来提炼刚性解析空间的定义(以便给形式概形配一个刚性解析的一般纤维),他便(二十三年后)记成了:我曾斩钉截铁地坚称,不需要簇的新概念,因为“我的”形式概形已够用一切。(记忆的失灵,倒常常把事情“办”得挺好……)然而,连 $K^*$(我的第二条主线)也并非来自形式概形。反正,这里依旧是那句老话:案子早已审理完毕!}——然而不过几天前,我还刚刚以世上最平静的心情写下几页,谈的正是Serre在某场\emph{葬礼}里将要扮演的某个关键角色、某个“支柱”的角色。唉,这一回可好,我正正地站在\emph{葬礼}之中,就在我鼻子跟前、电话线的另一头,而且恰恰就是这位Serre本人,一如既往地从容自在,分明怀着世上最好的信念!(话说回来,我怎么也想象不出Serre会不怀好意,尤其在谈到数学的时候……)

我自然没有心思争辩,Serre更没有,不过到底还是有一场东拉西扯的闲谈,五分钟或十分钟。这十分钟真是用得其所:让我亲身摩挲到一场\emph{葬礼}的可触真实——色、香、味,一应俱全;那场葬礼,因为我一味只盯着纸看,已经变得有点遥远了。

我不得不说出的头一件事,就是“刚性解析空间”这个名字本身,是我给起的——还暗示,如果我未曾明说的话,那是在我还独自梦想着这些我如此命名的东西的时候。Serre有点愣住——显然,他同样不记得;但同样明显,我不是在编故事。不过不要紧,名字终究只是名字,而且毕竟“自然得很”……这句“自然得很”明白地暗示:正因为太自然了,它已不再意味着什么;任何人只要鼻子对着那件东西,都免不了就叫它这个名:“刚性解析”。总而言之,我的朋友就这个名字,给了我一句他并无此意的恭维——却配着“就这点小事……!”的调门。再说,关于这个,我什么也没发表过,不是么?那就无话可说了……

我越来越瞠目。发表也罢,不发表也罢,对我而言,现实分毫不改。一个女人怀胎九月,把孩子生到世\dmark{}\origpage{1138}上,孩子已经又蹦又跳、好好的,有人对她说:这不是你的孩子,因为什么也没发表,而且她也没本事
出示出生证——对说这种话的人,她必定当面嗤笑。说实话,我没有冲着Serre笑——那不是我的
做派;何况我还没回过神来。我也没有想争辩:连Tate本人在他的笔记里,都对我在理论起步中扮演
的角色没有丝毫讳言(此事,Serre看来同Remmert忘得一样干净
\fnote{880}{(*)}{我再一次感到:“无论如何,案子早已审理完毕”。既然Tate说他“完全忠实地”遵从我提供的一位总施工,那好——不过是个总施工罢了,一张潦草的图,随便哪个小把戏都能在沙子上画出来;一股淡淡的Grothendieck味儿,没错——Tate还算客气,真是够朋友,居然费神提了一笔……});而1972年,当我写“主题概要”\fnote{881}{(*)}{即那份日期为1972年的文本,对我当时数学贡献的一份略嫌干巴巴(也不太启人)的概览,为申请法兰西公学院教席而写(该席位后来给了 J.~Tits)。这份文本,加上更周详的历史评注,将刊于《数学反思》第3卷。《引论》之三(“罗盘与行囊”)里谈到过它。“主题概要”5 e)中,我写道:\textbf{刚性解析空间。}受“Tate椭圆曲线”之例、以及完备离散赋值环上“形式几何”之需的启发,我达成了完备赋值域上刚性解析簇观念的部分表述;它在其后 J.~Tate对这一观念的首次系统研究里发挥了作用。此外,我在特征 $>0$ 域上的代数簇上引入的“晶体”,有时可以解释为特征零域上某些类型的刚性解析空间上带可积联络的向量丛;这预示着特征 $p>0$ 的晶体上同调、与特征零的刚性解析簇上局部系统上同调之间,存在着深刻的关系。}——这还是“如假包换的Grothendieck”!——其中影射此事时,Serre也未曾有过一丝
皱眉(他的记性,想必从那时起就开工了)。反正都是白费劲——分明只要没有发表,我说什么都会被
当作刮过耳边的风……

可“没发表”三个字让我心头一动,我顺势接下去:恰恰,我的著作的很大一部分,就是没有发表的东
西,口耳相传的东西。我感到Serre又是一愣——在他看来,这想法似乎有点荒诞,像一种用词上的自相
矛盾:“著作——没发表……”,在他那里,这两者仿佛凑不到一块儿\dmark{}\origpage{1139}。我说出了“Motif”这个词;他
立刻扑了上去:这一回,他要纠正我那套葬礼观念,兴冲冲地向我宣布:两三年前,恰好出了一整本关
于 Motif 的书——真的,在“Motif”这一章上,我没有什么可抱怨的!

“那么,那本出名的书,你亲手捧过吗?”我问他(问得正巧——我琢磨着问他这个有趣的问题,已有
一阵了)。

捧过——但我也许是想笑吧,Serre回敬道;他当然熟悉那本书;他甚至像通读过全书的人那样谈起
它;想必他确实读过。倘若里面没有什么特别的东西,我本可以免了这一问——明摆着没有;可是(人
就是这样,我毫无办法!)我还是问了!见他似乎没听懂问题的意思,我告诉他:去年我捧起它时,我
几乎不敢相信自己的眼睛。

\pmark 我不得不说出了“诈骗”这个词,但我感到,那是一个纯粹委婉的说法。照我真正的感受——写
这几行时依然如此——那是一场\textbf{无耻之行};可我忍住没有把它说出口。说到底,我心里明白,用哪个
词都无所谓:十五年来,一直“馊得厉害”,而Serre选择什么也感觉不到(那正是我几天前刚写下的),
无论我说什么,都不会“通行”。

Serre这下打开了话匣子,仿佛他等的正是这个。诈骗?你做梦吧,可怜的朋友;可这本书正是Deligne
本人写的,还写得漂漂亮亮呢,没错——行,人人都很清楚是你引入的 Motif,可这不构成每次一说
到“Motif”这个词就要重念一遍经的理由,对吧?更不必说你从来没发表过一行字,而你的瑜伽靠的是
未证明的猜想(我觉得听见另有某人借Serre的嘴对我说话……);而那本书的妙处恰恰是:它不
用任何猜想,实际上,它不用你当年做过的\textbf{任何}东西……

语气斩钉截铁、不容回嘴,出自一位深知自己在谈什么、再无可学之人——还带一丝恼火:一个有
点赶时间的人,被一个死心眼的大胖子缠着,偏\dmark{}\origpage{1140}不肯明白最明摆着的事情。那气氛,容不
得人对任何东西稍作停留——一切早已定案、拍卖成交,分明如此。Serre的公理——关于这一行的伦
理、何为要紧、何为次要——分明变了;对此无可奈何。我只能连他的新公理一起,照单全收。

于是我无可奈何,只在“猜想”上接了话!我本可以对他说:Weil猜想也是猜想——然而他也好,谁也
好,从未想过把它们不当回事——但确实,那些猜想,Weil费心发表了!可我恰好在写“第六根钉子”
(钉向我的棺材)
\fnote{882}{(*)}{指的是那组注(n\textsuperscript{o}s~176$_{1}$至176$_{7}$)——我正在给它做最后润色;在那组注里,我拆解的恰恰是围绕 Motif Galois群观念、围绕Galois--Poincaré-Grothendieck范畴(应景地被封为“塔纳卡”的)的骗局——一场由Deligne策动、(最初)由“棋子”Saavedra经手的骗局……},于是我把话岔到“Motif Galois群”上:它可没有任何
“猜想”的成分;我发展过一整套高度精确的理论,关于Galois--Poincaré型的范畴——那正是那本出
名的书里用到的基础观念之一,而书里似乎认为无须对我这个人做任何影射。

Serre又扑向了这个影射;这一回,他又能纠正我的葬礼幻觉了——那套理论,白纸黑字发表在一本
书里,而且还是我的一个学生写的:Saavedra
\fnote{883}{(*)}{即同一位Neantro~Saavedra~Rivano的那本出名的“塔纳卡范畴”(原文如此),《数学讲义》第265种(1972),Springer出版公司。}
——难道不正是我让他做的那篇论文吗?这一回,分明又是一本他烂熟于心的书;他不止一次引用过
它
\fnote{884}{(*)}{我还知道:每当Serre有机会引用这本(几乎)不提我名字的书——而他本人并不觉得有什么不对——他总(出于我不知道的什么顾忌)小心地同时提一提我这个人。肯费这番神的,他该是最后一个了……}。“那么,那本书里,就再没有什么让你吃惊的了”——我又问他(而这一回,分明我早已知道
答案)。

没有,显然不曾让他吃惊:我的名字不在书里出现,无论是作为主题的作者,还是作为其中引入的附
属观念(诸如 Motif、晶体等等诸如此类)的作者——那些观念\dmark{}\origpage{1141}从头出现、作为例子展开,出自我
的又一位学生之笔,依旧没有我的名字。而这一回,Serre看来并没有记忆失灵——他还记
得(至少眼下……)那些观念归谁。倘若我这位朋友确有“失灵”,那总归不在“记性”这一层……

我们又就“塔纳卡范畴”这个名字聊了几分钟:我暗示,我认为它是一场骗局;而Serre有根有据,觉
得它再贴切不过。这一点,其实早在挑起这只新野兔之前,我心里就明白;我也知道\textbf{为什么}这名字,
对那位曾是我的朋友的人,如此贴切——而我这个孕育并分娩了这件东西的人,却觉得有话要说。

一如我们之间的常态,是Serre斩断了谈话——事实上,谈话确实拖得够久了。自始至终没有一刻“沟
通”;所以它留给我那份不满与不谐,想来正因为此。然而,正像最近收到的那两三封短信——以更强
的、无可辩驳的力量——这场简短的谈话,教了我许多。都是“知道”的事,没错,却被打了对折;知道,
却不信!而那份挫折感(至今没有消散),正是我在抗拒接纳那个讯息的征兆。

一个不速的讯息,当然。几个月前,我还毫不怀疑:Serre(一如我对他鲜活的记忆——犀利优雅、方正
不苟的化身),当他(迟做总比不做好……)凭借阅读“收获与播种”这份天赐的文本,得知某场\emph{葬
礼}的种种丑行(他当然万万料不到,可怜的人……)之时,必将拍案而起,投身这场战斗
\fnote{885}{(*)}{写“这一回”时,我想到了另外两次:我两次下本钱,试图向那个出名的“数学共同体”传递讯息——甚至动员它。第一次是1970年,我告别数学舞台之际,为了科学建制与军事机器的串通。第二次,层次低些,只在法国同事之间,为了一项打击在法外国人的不公正条文。(关于这一点,见“我的告别——或:外国人”一节,n\textsuperscript{o}~24.)两次,我的努力都遇到普遍的冷漠;Serre并不比我刚离开的那个圈子里的任何别的朋友(Chevalley与Samuel这两位除外)更例外。人们不妨就“收获与播种”这块大砖头,在那个建制里——Serre首当其冲——将产生(或不产生)何种效应,打个赌……}。这幅埃皮纳勒版画式的想象,在\dmark{}\origpage{1142}最近几个星期里消散了——多亏一次平淡的书信往
还。而昨天,我有幸亲眼看清,毫无半点疑义:Serre置身\emph{葬礼}之中,早已既久且深,而且他在其中
颇为受用。这一点——无须说(我也绝无半点讽刺)——是抱着世界上最好的信念的!

其实,我早已明白:“善意”绝非我大半生所觉得的那样,是一件如此简单、如此非黑即白的事。有一
种最常见的“善意”,不过是给自己递烟幕:一面体面的旗帜,遮着有时可疑的货物。我们的心灵由层
层叠叠的膜构成;目光越细,越看见某一层的“善意”,如何有时给底下一层的骗局充当掩护与托辞。

至于Serre的善意,我继续这样信任他:他绝不会写一本书、在实质上使用他人的思想而不明白说出
——哪怕那些思想从未发表、除了传达给他的人(倘若还在世)与他本人无人知晓。也就是说,我相
信我知道:Serre绝不会写昨天我们谈到的那一类书。我甚至可以说:单是\textbf{像Serre、像我}这样的人\fnote{886}{(*)}{在这里说“Serre或我”,我实际想的是:五十年代我们同属的那个圈子里的任何一位成员——那个圈子,我在《愚妄与更新》第三、四部分里试图略加勾勒,尤其在“Bourbaki,或我的大机遇——及其反面”一节里。不过,确实,我发觉即便在这个小圈子里,我也知道有两位成员“变坏了”(《收获与播种》里在适当的地方谈到过)。},
写一份面向公众的文本(这里是数学文本),就牵动根深蒂固的职业良知反射;它们(我相信)会消除、
或至少纠正某些“记忆的失灵”——这类失灵,在昨天那样一场随意的闲谈里,本算不得什么
\fnote{887}{(*)}{因此我毫不怀疑:倘若Serre是一本讲刚性解析空间的书的作者或合著者(Remmert正是),他不会由着自己滑上那条“自然的坡道”——把那位理当被略过的人略过;他会越过那些对那条坡道颇为迁就的“记忆失灵”——而他在一次私人谈话里,恰恰乐得滑上那条坡道。同样确实的是:若在十五年前,凭我当时认识于他的那份严谨,他不会任自己滑上这样一条坡道——我以为如此,哪怕只是在一次私人谈话里……}。而这一切,正合三星期前\dmark{}\origpage{1143}我在“什么都不像的东西——或干涸”一注(n\textsuperscript{o}~171(v),部分(c))里写下的话的方向:“我知道,Serre同我一样,绝不会想到跟着狼嚎,在‘人人’抢、骗、拆的地方去抢、去骗、去拆。”

话虽如此,我发现:这一切并不妨碍Serre——至少在某些情形里——实实在在受用于\textbf{别人的}抢、
骗、拆;而且是公开的、明摆着的,“在大庭广众之间”、“在聚光灯下”。他当然可以“怀着最好的信念”
如此——他不必脏自己的手:他只消对别人的掠夺、操纵与拆毁毫无保留地祝福;更何况他分文不取:
他不夸耀他人劳动的果实,却觉得别人(我几乎想说:那些领了专营执照的经销商)当众玩这种把戏
挺好。他收下的“收益”,比(边上有点不干不净的)出版物、比另一些人垂涎的银行账户,更精微。可
见它们毕竟举足轻重,足以造成我所认识的那个人令人瞠目的变形:今天(说不出从何时起),他闭
着眼睛、捂着鼻子,参与着普遍的败坏\fnote{888}{(*)}{这一对参与集体败坏的确认,与“掠夺的空白授权——或‘高尚行当’”一注(n\textsuperscript{o}~171$_{4}$,其名与其所指相符)——里所作的确认(针对1980年3月某讨论班的听众)相合;尤其打字稿第1090页第二段。}。

\subsubsection*{e.\ 最后一刻——或一个禁忌的终结}
\tocsubsub{e.\ 最后一刻——或一个禁忌的终结}

(6月18日)昨天我犹豫过:要不要给“家庭相册”那条注(n\textsuperscript{o}~173)再添第四部分,好“趁热”报
道前天与Serre的那通电话。诚然,那通电话留给我的是一种“不满之感、不谐之感”(如我昨天所写)
——而且这还是委婉语:那份刺痛逼近焦虑。这份不安催促我回到那段插曲,像对待一颗已经熟透的
脓疮:早该切开引流了\dmark{}\origpage{1144}。还有惯常的迁延。几个星期以来,朗格多克科学技术大学的复制室,就等着人
把《收获与播种》那册出名的分册IV的续篇送去——它迟迟不肯分娩;眼下时间已紧,就得抢印抢装,
赶在学院7月15日年度闭馆之前,何况不只我一家:学年将尽,各色论文蜂拥而至,人家有优先权。总
之,我对自己说:一本书总得知道收尾;若我继续往里塞“最后一刻”,明年这时候也完不了;够久了
……

可我终究还是动笔了——哪怕《收获与播种》的印刷要等到开学!它已经等了十五年(三十年也差
不多),再等两三个月无妨;只求容我从容看一看我该看的,说一说我该说的,不被种种“截止日期”推
搡……

那是艰苦的一天劳作——不如说一夜加大半个上午——我要让这份“加料”文本赶上今天的邮班。
事已成矣。

此刻,我有一种走到某项\textbf{必做}之工尽头的感觉。我忽然觉得轻了,仿佛卸下了一副我背了不知
多久、却从未察觉的重负。那想必是一种顽固幻觉的重量;它早在四十年代末,就在我之内扎下根来:
当时,一种领来的身份在我之内绽放——某个“共同体”(数学的)成员的身份、某个圈子的身份;那个
圈子,在我眼里满盛着温暖与生命。这种新身份的绽放,我在《愚妄与更新》里——“受欢迎的陌生
人”与“‘数学共同体’:虚构与现实”两节(n\textsuperscript{o}s~9、10),还有“Bourbaki,或我的大机遇——及其
反面”(n\textsuperscript{o}~22)一节——谈到过。诚然,这份认同,被1970年我离去前后的事件一扫而空,无影无
踪——那是在我投入战斗行动之后。可是,如今回顾,我意识到:对那个我已离弃、不再自认同属的圈
子,毕竟还系着一条\textbf{纽带}——一条也许无形、却力量巨大的纽带;它属于那份“过去的重量”(去年,
同名一节,n\textsuperscript{o}~50,我开始瞥见它)。我离开了那个圈子,没有回头的念头;可那个“家”\dmark{}\origpage{1145}——我为了
另一场历险而离弃的家——的一幅\textbf{图像},却在我之内活着,维系着那条纽带。这
幅图像,想必自我离去(以及更早,当然)直到《收获与播种》的思考,大体静止。《收获与播种》开始
修饰我对某段过去的图像,并勉力把现下的元素——常常令人错愕、常常不请自来——纳入其中。我
最终认了那份明摆着的事实:在我曾自认同属的圈子的后续、以及(据说)整个数学世界里,人心与
风俗发生了一场惊人的\textbf{败坏}。我意识到:这场败坏并非始于昨日;早在我离去之前,我就有足够的
时间,贡献过我的一份(至少,《愚妄与更新》的思考里瞥见过一份)。可是我有印象:我离去之
后,这场败坏有一场疯狂的升级;我的一些前学生,在其中扮演了头等的催化剂角色。

不管怎样——在对\emph{葬礼}的调查接踵而至的层层揭示之中,我心里始终维持着一种心照不宣
的“禁忌”,围着我旧日朋友中那些属于我青年时代那个圈子的人。我简直无法设想:他们中会有谁,
被这场我不断确认的深层败坏严重侵蚀、或“伤到”。当我说起“教团全体”对那些(至少在我眼里)超
出想象的行动的纵容时,我之内想必设有一道内在的“除外条款”,把那些在我看来应“立于一切嫌疑
之上”的人排除在外。他们分明毫不知情——他们想必在忙别的,肯定如此——可别怪他们!大致是这
样的调门。至于我的长者中最年长的几位,这种看法,我愿意相信,符合真实——或至少符合真实的
一个侧面。但对Serre、Cartier、Borel、Tate、Kuiper、Tits以及我熟知的其他人,却断然不然:他们与我
同代,正当盛年,完全融入我此刻检视的这个圈子,直到今天仍在其中行使着不小的权力、定着调
门,正如某些新来乍到者——他们在长者们的祝福之下,终于在那个圈子里结成了一伙毫无顾忌的
“黑手党”。

于是,在我据以看现实的那幅图像里——现实透过\emph{葬礼}这台头等的“显影剂”呈现——有一
道顽固而触目的矛盾。正是这道矛盾,无疑,\dmark{}\origpage{1146}在某个层面上被知觉、又在另一个层面上被否弃,造成了方才
说的那种逼近焦虑的“不安”——那焦虑,泄露了一场\textbf{分裂}。而那个人——比任何人都更能在我眼中体现这个圈子,体现这群被我内心的某个角色执意视为“亲近者”的人;而他本人,也曾是他们之中“最亲近”的一个——就是Jean-Pierre~Serre。正因此,被回避的矛盾之结,主要系在他身上。

\pmark 我开始怯怯地触及这道矛盾,不过六星期前,在“无用的细节”(n\textsuperscript{o}~171(v))一注的第一
部分(5月4日)。这场思考在同一注的第三部分(5月27日,即三周之后)——“什么都不像的东西——或
干涸”——里大为深化。一周之后(6月11日),我顶着顽强的内心抗拒,在本注的 c.部分(“众人之上的那
一位——或首肯”)里,重新回到Serre其人。这一回,Serre在\emph{葬礼}中的关键角色,终于大白于光
天之下。这是我对\emph{葬礼}的理解中的新的一大步——可矛盾的结,依旧没有触及!Serre其人,在我
这里(仿佛什么也不曾发生过)依然是无畏无瑕的“优雅”与“方正”的化身。“禁忌”完好无损!

是前天的那个电话,炸开了这道矛盾,把我的鼻子“直接摁进”(\emph{葬礼})里——不管我乐意与否。
随之,当然,是强大的抵抗之力(方才提到过的)的立即动员,以维持现状,而不是承认矛盾:以这样那
样的方式记录在案,由此解决它。我可以这样做,也可以不这样做。

\pmark 我跨出了那一步——而且我为此高兴。报偿是即时的:一场\textbf{解放},表现为轻快与如释重
负之感;一场内在张力的松解,诚然,但更是一场重负的卸除。

在《收获与播种》里,唯一可与之相比的解放时刻,同样标志了思考的第一次大转折:那是《愚妄与
更新》里的“竞技数学”,紧随其后的是“旋转木马收摊了!”(n\textsuperscript{o}s~40、41)。我还有印象:我方才“跨出”的
这新的一步,接续着去年那一步。何以如此\dmark{}\origpage{1147}、在何种意义上如此,我一时说不好。当时那声凯旋的欢呼
“旋转木马收摊了!”,肯定是过早了(次月我就察觉了)。可我刚迈的这新的一步,至少是又一步,把我带离那
台出名的旋转木马更远一步。将来会告诉我,实际情形是否如此。

经过昨天与6月11日的思考,我觉得已抵达一幅对\emph{葬礼}不那么含混的图景。此前,尤其
是这个“第三平面”模糊不清。6月11日的思考,让它“以可触的方式”化身于Serre其人;而这具化身,又
在昨天的思考里,取得了(这样说毫不夸张)一切具体的轮廓。

\pmark 说到底,在《收获与播种》这整个第四部分里,关于我与Serre关系的思考,在我看来,是对我理
解\emph{葬礼}最关键的部分——在“侦查补遗”与那幅数学大都会阴沟的浓墨重彩的图景之外。诚然,
倘若我不曾出于对我交给自己的题目的敬意,以我所能的全部细心,硬着头皮做这份“侦查的归置”,
并且尽力照亮沿途出现的每一个稍暗的角落,关于Serre的这场思考,恐怕也不会问世;而我对\emph{葬
礼}(以及我自身对它的牵连)的理解,将依旧如前含混。研究的工作里,一切互相维系!

\pmark \emph{葬礼}最后部分里那场思考的最有分量的部分,实际上是“最后一刻”才出现的。照理,这一
部分的“句号”,两个半月前(4月7日)就已画下。当时只剩十来页要重抄清稿、补几条脚注(一年前五
月末,情形也一样……)。意外之事,随即接踵而至:Zoghman来访——来读这(原则上已完成的)最后一
部分、向我提意见。它们凝结成了大约三百页的追加文字——其中,就有这几页:在葬礼(至今被回避)
的光照之下,我重新审视我与Serre的关系。

